\chapter{电源、时钟与信号完整性}

前面八章默认了两件事：电源永远是稳定的电压，时钟永远是干净的方波。真实机器里这两条都要靠工程手段挣出来。电压会因为瞬态电流而跌落，时钟会因为走线长度不一而偏斜，信号边沿会在走线上来回反射。这些问题不写在真值表里，却决定真值表能不能成立。本章把这三类问题从"经验之谈"变成可计算、可测量的对象：每个现象先给物理原因，再给数量级估算，最后给测量和排查方法。

\begin{longtable}{L{0.18\linewidth}L{0.4\linewidth}L{0.32\linewidth}}
\toprule
对象 & 失效时的样子 & 本章回答的问题 \\
\midrule
电源分配 & 瞬态电流让电压跌落，逻辑误翻转 & 去耦网络怎样按时间尺度分层 \\
时钟分配 & 各触发器时钟沿不齐，时序预算作废 & 偏斜和抖动从哪里来、怎么限 \\
信号完整性 & 边沿振铃、过冲、回沟、串扰 & 什么时候必须把走线当传输线 \\
\bottomrule
\end{longtable}

\begin{keyidea}
电源、时钟和信号共享同一个物理视角：它们都是沿导体传播的电磁过程，寄生电阻、电容和电感在低速时可以忽略，在边沿变快、电流变大时必然现身。本章的每条规则都能还原成这三个寄生参数的数量级比较。
\end{keyidea}

\section{一、电源分配：稳压、去耦与电源完整性}

本节从"电源"这两个字背后的一整串工程问题开始。数据手册把
供电写成一行参数：VCC = 5V $\pm$5\%。真实板子上，这一行要
靠三类手段兑现：稳压器把未经稳压的输入变成标称电压，去耦
网络对付稳压器来不及响应的快变化，电源与地平面的几何形状
决定噪声在板内怎样传播。三者缺一，逻辑完全正确的电路也会
偶发失误。

本节按时间尺度从慢到快组织：稳压器管 ms 级以上的平均供电；
瞬态电流与电压跌落说明为什么稳压器天然管不了 ns 级事件；
去耦电容按时间尺度分层，接管从 ns 到 ms 的全部空档；目标
阻抗把这些层次统一成一条随频率变化的阻抗曲线；最后，地平
面与回流路径回答"电流回哪里去"这个同样关键的问题。

\topic{线性稳压与开关稳压的取舍}两类稳压器的差别在一个物
理事实上：多余的那部分电压去了哪里。线性稳压器（78xx 系
列、各类 LDO）内部是一只工作在线性区的调整管，输入高出输
出的电压直接降在管子上，以热的形式耗散。效率因此几乎只由
电压比决定：$\eta \approx V_{out}/V_{in}$。5V 降到 3.3V，
效率上限就是 $3.3/5 = 66\%$，与器件好坏无关——这是物理限
制，不是工艺限制。开关稳压器（LM2596 一类 buck）的做法完
全不同：调整管只有全开、全关两个状态，占空比决定平均输出，
电感在每个开关周期里储能再放能，把能量一批批搬到输出端。
理想开关不耗能，实际效率约 75\%--90\%，损失主要在开关过程、
电感电阻和续流二极管上。

给 5V 转 3.3V、1A 这一路算一笔账。用 LDO：输入功率 5W，
输出 3.3W，差值 $(5-3.3)\,\mathrm{V} \times 1\,\mathrm{A}
= 1.7\,\mathrm{W}$ 全部变成热。1.7W 对 TO-220 封装不是小
数目：无散热片时结到环境热阻约 50 $^{\circ}$C/W，温升约
85 $^{\circ}$C，室温稍高结温就逼近上限，散热片省不掉。换
LM2596 按效率 80\% 估：输入只需 $3.3/0.8 \approx 4.1$W，
损耗约 0.8W，不到线性方案的一半。压差越大差距越悬殊：输
入换成 12V，线性方案效率只剩 $5/12 \approx 42\%$，开关方
案仍在 80\% 附近——同一负载，发热相差五倍以上。

代价在纹波和外围器件。LDO 的输出接近一条直线，纹波以
$\mu$V 到 mV 计，适合给模拟前端和噪声敏感电路供电。开关
稳压器的输出叠加着开关频率的锯齿：LM2596 工作在约 150kHz，
输出纹波典型几十 mV，还混有开关瞬间的尖峰；给数字逻辑供
电通常无妨，给 ADC 参考电压或射频电路供电就要追加滤波。
外围器件也更多：buck 需要电感、续流二极管（或同步开关管）
和输入输出电容，电感的体积和成本常常超过芯片本身。7805
还要求输入至少高出输出约 2V——这个最小压差称为 dropout，
5V 输出就至少要 7V 输入；现代 LDO 可低至零点几 V，"低压
差"三个字正是这么来的。选型因此是一道算术题而不是口味题：
压差小、电流小、噪声敏感，用 LDO；压差大、电流大、效率优
先，用开关。

开关与线性还可以串联使用，这正是经典组合"开关预稳压加
LDO 后级滤波"的理由。LDO 的反馈环路不仅稳压，也抑制输入
纹波：数据手册的电源抑制比（PSRR）在低频段常有 60dB，输
入 100mV 的纹波到输出只剩 0.1mV；但 PSRR 随频率升高而下
降，150kHz 的开关纹波还剩多少抑制，必须查曲线而不是猜。
正确用法是让 buck 把 12V 高效降到 4V 上下，再由 LDO 滤到
干净的 3.3V——效率由开关挣，干净由 LDO 挣。直接拿开关输
出给噪声敏感电路供电，等于只算了这笔账的一半。

精确一点，线性稳压器的效率应写成 $\eta = V_{out} I_{out}
/ (V_{in} (I_{out} + I_Q))$，$I_Q$ 是器件自身的静态电流，
7805 约几 mA。重载时 $I_Q$ 可忽略，效率就是电压比；轻载
时 $I_Q$ 占比上升，效率比电压比还低。开关稳压器轻载时也
要为每个开关周期付出固定损耗。所以"效率 80\%"是额定负载
附近的数字，空载时两类稳压器都不好看——给电池供电的常开
电路选型时，数据手册的静态电流一栏比效率曲线更关键。

测量上，两类稳压器一眼可分。量输入、输出电流：线性稳压器
的输入电流几乎等于输出电流（只差几 mA 静态电流）；开关稳
压器的输入电流按功率折算明显更小，12V 转 5V/1A 的 buck 输
入电流只有约 0.5A。器件烫手时，先算 $(V_{in}-V_{out})
\times I_{out}$ 是否超出封装散热能力；纹波异常时，示波器
AC 耦合看输出：近直线是线性，规则的 kHz 锯齿是开关，锯齿
幅度超标多半是输出电容老化失效。

\begin{longtable}{L{0.2\linewidth}L{0.34\linewidth}L{0.34\linewidth}}
\toprule
关键问题 & 线性稳压（7805 / LDO） & 开关稳压（LM2596 buck） \\
\midrule
多余电压的去向 & 降在调整管上变成热 & 开关搬运，理想不耗能 \\
效率（5V 转 3.3V/1A） & 66\%，浪费 1.7W & 约 80\%，浪费约 0.8W \\
输出纹波 & $\mu$V--mV 级，接近直线 & 几十 mV，含 150kHz 锯齿 \\
必需外围器件 & 输入输出电容各一只 & 电感、二极管、输入输出电容 \\
输入压差要求 & 7805 约 2V，LDO 可更低 & 输入范围宽，大压差也能工作 \\
适用场合 & 小电流、噪声敏感、压差小 & 大电流、压差大、效率优先 \\
\bottomrule
\end{longtable}

\topic{瞬态电流让电压跌落}芯片不是按平均值耗电的。CMOS 门
只在翻转瞬间从电源取电流：输出由低充高，电流给负载电容充
电；由高放低，电荷泄到地。一组门同时翻转——32 位总线换
数、计数器多位同翻、时钟沿后整片组合逻辑一起动作——尖峰
叠加成一个 ns 级电流脉冲。教学规模就容易达到 0.4A 在 5ns
内出现又消失的量级（16 路输出各取 25mA 即 0.4A）；大型芯
片的同时开关电流以 A/ns 计。

电流尖峰本身不可怕，可怕的是它流过寄生电感。任何一段导体
都有电感，经验法则：每 mm 走线约 1nH，2cm 电源走线就是约
20nH。电感的本性是反对电流变化，感应电压 $v = L \cdot
di/dt$。代入：$di/dt = 0.4\,\mathrm{A}/5\,\mathrm{ns} =
0.08\,\mathrm{A/ns}$，$v = 20\,\mathrm{nH} \times
0.08\,\mathrm{A/ns} = 1.6\,\mathrm{V}$。尖峰出现的几 ns 内，
芯片实际看到的电压从 5V 跌到 3.4V——而 5V $\pm$5\% 只允
许 $\pm$0.25V，这一次跌落是允许值的六倍多。后果不是"信号
差一点"这种模拟意义上的劣化，而是门延迟突变、存储节点电
荷平衡被破坏，触发器采错或保持不住，表现为与数据模式相关
的偶发错误。

指望稳压器纠正这种跌落，是错配了时间尺度。稳压器靠反馈环
路工作：采样输出、比较基准、调整管子，环路带宽折算成响应
时间是 $\mu$s 级——LDO 几 $\mu$s，开关稳压器还受开关周
期限制，150kHz 一个周期 6.7$\mu$s。电压跌落发生在 5ns 里，
比稳压器的反应快三个数量级，稳压器连"出事"都来不及知道。
ns 级事件的供电只能靠储能就在芯片身边的元件：电容。电容
与芯片之间的路径电感决定最初几 ns 的电压缺口，容值决定之
后能撑多久。

电容要出多少货，可以算出来。0.4A 持续 5ns，电荷 $Q = I
\times t = 2\,\mathrm{nC}$。允许电压因此下降 0.1V，需要
$C = Q/\Delta V = 20\,\mathrm{nF}$——一只 100nF 在忽略寄
生时绰绰有余。但前提写在上一段：电容到电源脚的路径电感必
须远小于 20nH，否则最初几 ns 仍是 1V 量级的缺口。这就是
"就近"排在"够大"前面的原因：贴着电源脚的 100nF，远胜板
子另一端的 10$\mu$F。

大型处理器把同一个方程推到极端：核心电压 1V 上下、瞬时电
流变化几十 A、边沿 ps 级，di/dt 以几十 A/ns 计——哪怕
1nH，也是几十 mV 起步的坑。处理办法仍是同一套时间尺度分
工，只是层数更多：板上电容、封装电容、片上电容各管一段，
电源引脚数以百计。芯片封装内部也要算这笔账：每个电源脚自
带几 nH 键合线电感，多个输出缓冲同时翻转时，片内电源轨和
地轨跟着晃动，称为同时开关噪声（SSN），§四会从信号串扰的
角度再看它一次。教学机的 0.4A/5ns 与处理器的几十 A/ns 相
差六个数量级，方程是同一个。

排查时把示波器探头直接搭在芯片电源脚与就近地之间，地线越
短越好——长地线自己就是电感，§五会专门讲。用单次触发捕
捉总线翻转瞬间的电源波形。错误指纹也值得记住：跑固定数据
死、换一组数据就过，这种与数据模式相关的偶发错误，第一嫌
疑就是电压跌落。

\topic{去耦电容按时间尺度分层}没有一只电容能覆盖全部时间
尺度，原因仍是寄生：每只实际电容都串联着引脚与内部结构的
电感（ESL）和电阻（ESR）。大容值的电解、钽电容体积大、ESL
大，对快事件相当于不在场；小容值的陶瓷电容 ESL 小、响应
快，但储能少、撑不久。去耦网络因此按时间尺度分层，每层只
管自己来得及管的事：

\begin{longtable}{L{0.16\linewidth}L{0.22\linewidth}L{0.24\linewidth}L{0.26\linewidth}}
\toprule
层级 & 典型容值与介质 & 服务的时间尺度 & 位置规则 \\
\midrule
芯片旁 & 100nF 陶瓷 & ns 级尖峰 & 贴电源脚，路径几 mm 内 \\
局部 & 1--10$\mu$F 陶瓷/钽 & 几十 ns 的成组翻转 & 每个功能模块一只 \\
板级 & 10--100$\mu$F 电解 & $\mu$s--ms 负载变化 & 电源入口，撑到稳压器响应 \\
稳压器输出 & 数据手册指定 & 环路响应以内 & 紧贴稳压器，关系环路稳定 \\
\bottomrule
\end{longtable}

距离规则逐层放宽的理由：层级越高，事件的 di/dt 越小，路径
电感惹出的 $v = L \cdot di/dt$ 越小，允许的距离就越大。ns
级尖峰只给几 mm 的往返路径，100nF 必须贴在电源脚旁；$\mu$s
级事件的 di/dt 小了三个数量级，10$\mu$F 放在 cm 级距离照
样有效；板级电解甚至可以只管电源入口。容值规则逐层加大约
百倍的理由：每层既要自己存货够用，又要能在下一层耗尽之前
接力。100nF 在 0.1V 压降下放完的储能是 10nC，约接五次 2nC
尖峰；尖峰成群出现时，由旁边的 1--10$\mu$F 给它补仓，依此
类推，直到稳压器环路接管最慢的部分。

把这张表竖起来画，去耦网络的分层结构就更直观：越靠近芯片，服务的事件越快、电容越小、贴得越近。

\begin{pdfdiagram}
  \node[hw compute, minimum width=4.6cm] (vreg) at (0,4.4) {稳压器（输出电容按手册）};
  \node[hw memory, minimum width=4.6cm] (bulk) at (0,3.2) {板级 10--100$\mu$F 电解};
  \node[hw memory, minimum width=4.6cm] (loc) at (0,2.0) {局部 1--10$\mu$F 陶瓷/钽};
  \node[hw memory, minimum width=4.6cm] (byp) at (0,0.8) {芯片旁 100nF 陶瓷};
  \node[hw state, minimum width=4.6cm] (chip) at (0,-0.4) {芯片电源脚};
  \draw[hw bus] (vreg.south) -- (bulk.north);
  \draw[hw bus] (bulk.south) -- (loc.north);
  \draw[hw bus] (loc.south) -- (byp.north);
  \draw[hw bus] (byp.south) -- (chip.north);
  \node[hw label, anchor=west] at (2.5,4.4) {ms 以上，环路响应以内};
  \node[hw label, anchor=west] at (2.5,3.2) {$\mu$s--ms 负载变化};
  \node[hw label, anchor=west] at (2.5,2.0) {几十 ns 成组翻转};
  \node[hw label, anchor=west] at (2.5,0.8) {ns 级尖峰};
  \node[hw label, anchor=west] at (2.5,-0.4) {瞬态电流由此取用};
  \draw[hw return] (-3.6,4.2) -- (-3.6,-0.2);
  \node[hw label, rotate=90] at (-3.95,2.0) {事件越快，电容越小、越贴近芯片};
\end{pdfdiagram}
\diagramnote{去耦网络按时间尺度分层。储能自上而下逐层减小、响应逐层变快：每层只管自己来得及管的事件，耗尽之前由下一层补仓，最慢的部分留给稳压器环路。}

缺层的后果各不相同，指纹却是同一张：与数据模式相关的偶发
错误。只有 100$\mu$F 没有 100nF，快事件没人接，电源波形上
全是窄尖峰；只有 100nF 没有 $\mu$F 级，成群翻转时 100nF
瞬间耗尽，电压缓跌几 $\mu$s；100nF 装了却离电源脚 5cm，
路径电感把它与快事件隔开，等于没装——手工板上最常见的一
种。分层规则浓缩成三句话：小电容给快事件，大电容给慢事件；
越小的电容越贴近芯片；相邻层级容值相差约百倍，让每层都接
得住下一层。

手工板与面包板上的做法是对这张表的直译：每片芯片的电源脚
旁一只 100nF 陶瓷，引脚剪到最短、跨在电源与地之间；每块
面包板的电源入口一只 10$\mu$F；整板进线处一只 47--100
$\mu$F 电解。这套规矩从 TTL 时代沿用至今，器件换了几代，
时间尺度的物理没有变。最常见的错误是把电容插在离芯片两三
排孔之外"意思一下"——对 ns 级事件，那几 mm 的孔距就是电
感，电容约等于没插。

去耦是否到位，可以用对比实验量出来。示波器探头搭在某片芯
片的电源脚上，先记录当前纹波；再在同一个电源脚上并一只
100nF，纹波中的窄尖峰应明显变矮；把这只电容挪到几 cm 外，
尖峰又涨回来——同一只电容，距离不同，效果立见。这个实验
值得在面包板上亲手做一遍：它把"就近"两个字从规则变成亲
眼所见。

\topic{PDN 目标阻抗：把电压波动写成欧姆问题}前面把时间轴
切成几段、每段配一只电容，工程上还需要一个把所有段落统一
起来的设计目标。办法是换到频率轴：把"电压不许波动超过
$\Delta V$"与"负载电流阶跃最大 $\Delta I$"相除，得到一
条对全频率成立的阻抗上限——电源分配网络（PDN）的目标阻
抗 $Z = \Delta V / \Delta I$。5V $\pm$5\% 意味着 $\Delta
V = 0.25$V，瞬态 1A 意味着 $\Delta I = 1$A，目标阻抗
0.25$\Omega$（250m$\Omega$）。窗口一收紧数字立刻变小：
$\pm$0.5\%（25mV）对 1A，目标就是 25m$\Omega$；3.3V
$\pm$5\%（165mV）对 5A 瞬态，约 33m$\Omega$；1V 核电压、
几十 A 瞬态的处理器，目标阻抗以 m$\Omega$ 计。电源设计的
全部手段，都是让 PDN 阻抗曲线在整个频率轴上位于这条目标
线以下。

"足够多、足够近、足够多种容值"这条口诀，在阻抗曲线上看
得最清楚。单只电容的阻抗随频率呈 V 形：低频端是 $1/(2\pi
fC)$，随频率升高而下降；高频端被 ESL 接管，$2\pi f L$ 随
频率升高而上升；谷底是自谐振频率 $1/(2\pi\sqrt{LC})$。给
100nF 配 5nH ESL，自谐振约 7MHz，高于此频率它是一只电感；
10$\mu$F 配同样 5nH，自谐振约 0.7MHz。两条 V 形错开正好
接力：10$\mu$F 管 MHz 以下，100nF 管 10MHz 附近——这就是
"足够多种容值"。同值电容并联 N 只，阻抗近似除以 N——这
就是"足够多"。每只电容到芯片的路径电感直接抬高 V 形右半
边，距离翻倍、高频段阻抗跟着翻倍——这就是"足够近"。更
高频段（百 MHz 以上）分立电容无能为力，靠电源-地平面对的
平面电容和片上电容收尾；教学机到不了这个频段，规则却是同
一条。

目标阻抗可以直接用来"点兵"。仍取 5V $\pm$5\%、瞬态 1A，
目标 0.25$\Omega$。看 10MHz 这一点：一只 100nF 配 5nH
ESL，已过自谐振，阻抗约 $2\pi f L = 2\pi \times 10^{7}
\times 5 \times 10^{-9} \approx 0.31\,\Omega$，单只不达标；
4 只并联，阻抗约 0.08$\Omega$，达标。同一算法对每个关心
的频率点重复一遍，电容的数量与搭配就定了——去耦设计至此
从经验口诀变成算术。

目标阻抗也把测量变成明确的事：测出 PDN 阻抗随频率的曲线，
凡是冒过目标线的频段，就是偶发错误的嫌疑频段。排查顺序：
先算目标阻抗，再看曲线在哪个频率冒头，代入自谐振公式，就
能锁定该补哪一层。

把 10$\mu$F 与 100nF 两条 V 形曲线连同目标线画在同一张图上，这层"接力"关系就一目了然。

\begin{pdfdiagram}
  % 坐标轴（双对数示意）
  \draw[line width=0.55pt, -{Stealth}] (0,0) -- (9.5,0);
  \draw[line width=0.55pt, -{Stealth}] (0,0) -- (0,3.6);
  \foreach \x/\t in {0/10k, 1.8/100k, 3.6/1M, 5.4/10M, 7.2/100M, 9.0/1G} {
    \draw[line width=0.45pt] (\x,0) -- (\x,-0.09);
    \node[hw label, anchor=north] at (\x,-0.13) {\t};
  }
  \foreach \y/\t in {0/0.01, 1/0.1, 2/1, 3/10} {
    \draw[line width=0.45pt] (0,\y) -- (-0.09,\y);
    \node[hw label, anchor=east] at (-0.13,\y) {\t};
  }
  \node[hw label, anchor=north] at (9.0,-0.62) {频率/Hz（对数）};
  \node[hw label, anchor=east] at (-0.13,3.5) {$|Z|$/$\Omega$（对数）};
  % 目标阻抗线
  \draw[BookRed, dashed, line width=0.7pt] (0,1.4) -- (9.0,1.4);
  \node[hw label, text=BookRed, anchor=south east] at (8.95,1.45) {目标阻抗 250m$\Omega$};
  % 10µF 的 V 形：自谐振约 0.7MHz
  \draw[BookTeal, line width=0.9pt] (0,2.2) -- (3.32,0.36) -- (8.4,3.2);
  \fill[BookTeal] (3.32,0.36) circle (0.055);
  \node[hw label, text=BookTeal] at (1.1,0.62) {10$\mu$F\\自谐振 0.7MHz};
  % 100nF 的 V 形：自谐振约 7MHz
  \draw[BookAccent, line width=0.9pt] (2.17,3.2) -- (5.13,0.36) -- (9.0,2.5);
  \fill[BookAccent] (5.13,0.36) circle (0.055);
  \node[hw label, text=BookAccent] at (6.55,0.42) {100nF\\自谐振 7MHz};
\end{pdfdiagram}
\diagramnote{PDN 阻抗随频率的分布（双对数示意）。单只电容的阻抗呈 V 形：低频端按 $1/(2\pi fC)$ 下降，自谐振点触底，高频端被 ESL 接管按 $2\pi fL$ 上升。两种容值的 V 形错开接力，设计目标是让合成阻抗在全频段压在目标线以下。}

\topic{地平面与回流路径}去耦回答电流从哪里来，同样重要却
更常被忽略的是：电流回哪里去。每个信号都携带一个回流——
驱动器推出电流，沿信号线到达接收端，再沿参考面流回驱动器
的地。直流回流走电阻最小的路径，在整个地平面上铺开；高频
回流走阻抗最小的路径，而高频下阻抗以电感为主，电感最小意
味着环路面积最小——于是高频回流贴着信号线正下方的地平面
走。这是电磁场自己选的路：信号线与正下方的回流构成一对紧
耦合导体，环路面积被限制在"走线长度乘介质厚度"的窄条里，
介质厚度只有零点几 mm。

一旦在地平面上开缝，这条窄条就被撕开。信号线跨越分割缝时，
正下方没有铜，回流只能绕到最近的桥接点再绕回来，环路面积
从 mm$^2$ 级涨到 cm$^2$ 级。环路电感大致随面积增长，几十
nH 很平常；再代入本节开头那个 0.4A/5ns 的尖峰，又是 1V 量
级的噪声——只是这次它出现在地参考上，称为地弹：所有以地
为参考的信号跟着一起晃。大环路还是高效天线，§四讲电磁兼
容时还会回到这一笔。

分割不是原罪，跨缝才是。合理的分割把噪声悬殊的电路隔开：
继电器、电机驱动与逻辑电路分两地，唯一连接做在电源入口的
单点；模拟小信号区与数字区分地，唯一连接做在 ADC 正下方。
分割成立的判据只有一条：没有任何信号线跨越分割缝，所有跨
区信号都经过桥接点同行。信号密集交错、说不清哪根线会跨缝
的板子，宁可保留完整地平面、靠器件分区摆放来隔离，也不要
开一条逼着回流绕路的缝。

没有完整地平面时的折中方案，同样是把回流"贴着信号走"人
工维持出来。双面板可以织网格地：顶层地线横走、底层地线纵
走，过孔在交叉点相连；关键信号（时钟、复位）旁边再有意识
地并行走一根地线。面包板连网格也没有，那就把时钟线与其回
流地线绞在一起，把环路面积拧到最小。低速时绕路代价小，高
速时绕路直接报废时序与波形；但回流不是自动存在的，是被设
计出来的——这一点高低速相同。

完整地平面还有第二个身份：它是信号走线特性阻抗的参考面。
走线阻抗由导线与参考面的几何关系决定，参考面连续，阻抗才
连续；回流贴着信号走与阻抗连续，是同一块铜的两个作用。§三
讲传输线时会发现，阻抗突变处的反射与回流绕路处的噪声，追
到底常常是同一条地缝。地平面的完整性因此同时是电源完整性
与信号完整性两章内容的共同前提。

排查接地问题的手法有两个。一是查图：沿每根高速信号走线看
正下方的地平面是否连续，跨缝处就是第一嫌疑人。二是量地弹：
示波器探头的信号端与地线分别接同一地网络上的两个点，测到
的就是两点间的电压差，即回流惹出的噪声；偶发错误若与某几
根信号的活动强相关，先量这些信号的回流路径。

\begin{longtable}{L{0.28\linewidth}L{0.3\linewidth}L{0.3\linewidth}}
\toprule
现象 & 先怀疑什么 & 检查点 \\
\midrule
负载加重就复位 & 板级储能不足、稳压器限流 & 电源入口有无 $\mu$s 级缓跌 \\
错误与数据模式相关 & 100nF 缺失或离芯片太远 & 芯片电源脚上的窄尖峰 \\
某模块一动作全板出错 & 回流绕路引起地弹 & 该模块信号是否跨越地缝 \\
输出有规则 kHz 锯齿 & 开关纹波，判断是否超标 & AC 耦合测峰峰值对照容忍度 \\
LDO 烫手 & $(V_{in}-V_{out}) \times I_{out}$ 热耗 & 散热片与封装热阻 \\
\bottomrule
\end{longtable}

这张排查表的使用方式是从指纹倒推元件：先确认错误与数据模
式、负载、某模块动作中的哪一项相关，再按时间尺度对号入座
——ns 级尖峰找 100nF 的位置，$\mu$s 级缓跌找板级储能与
稳压器，地与信号一起晃找回流路径。电源问题极少需要换芯片，
答案几乎都在电容、距离与路径这三件事上。

\begin{keyidea}
电源完整性可浓缩为一条主线：稳压器负责 ms 级以上的平均电
压，去耦网络负责让 PDN 阻抗在所有频率下低于 $\Delta V /
\Delta I$，地平面的完整性决定回流是否绕路。按时间尺度分工、
按频率轴核算，是电源设计的全部方法。
\end{keyidea}

\section{二、时钟树：分配、偏斜与抖动}

第二章建立时序预算时，默认每个触发器在同一个瞬间看到时钟
沿。物理上这个瞬间不存在：时钟从源出发，沿走线传播，经缓
冲器放大，到达各触发器的时刻必然有先有后——这是\hwkey{偏
斜}（skew），确定性的、由结构决定的时间差。即使结构完全
对称，沿的位置还会被电源噪声、串扰和器件噪声推着随机游移
——这是\hwkey{抖动}（jitter），统计性的、由噪声决定的时
间差。

把两者分清有一个实际理由：修理手法互不通用。偏斜靠改结构
（等长、对称、补齐级数），抖动靠改噪声（去耦、屏蔽、换更
好的源），锁定问题靠改环路（带宽、参考、配置次序）。故障
定位的第一步永远是先问：这个时间差是确定的，还是随机的？
本节按"先结构后统计"组织：先看分配结构怎样把偏斜收小，
再把偏斜与抖动写回时序预算，最后看管时钟源头的两个部件：
锁相环与晶体振荡器。

\topic{时钟分配用等长的树}让所有触发器的时钟沿对齐，等价
于让"时钟源到每个触发器"的每条路径延迟相等。延迟来自两
部分：走线传播延迟与缓冲器延迟。FR4 板材上信号传播速度约
15cm/ns（光速除以介电常数的平方根，$\sqrt{4} = 2$），1cm
走线约 67ps；一级缓冲器几 ns。两种极端结构的对照最能看清
问题。菊花链：时钟先到第一个触发器区，再串到下一个，像击
鼓传花——8 级链、每级走线加缓冲 6ns，末端比首端晚 48ns，
而 50MHz 的周期一共只有 20ns。等长树：从根分杈，每级扇出
经缓冲器驱动下一级，各分支走线等长、缓冲级数相同，所有叶
节点的沿在结构上同时到达。

把数字代进树里。某设计用两级缓冲，每级 5ns。只要两个分支
的缓冲级数相同，$2 \times 5$ns 在两边相消，偏斜只剩走线
差：分支长度差 10cm，$10\,\mathrm{cm} \div 15\,\mathrm{cm/ns}
\approx 0.67$ns，这就是两个叶节点之间的偏斜。再看级数不
相等的情形：一个分支两级、另一个三级，仅缓冲延迟就差出
5ns，加上走线差共 5.67ns——级数匹配比长度匹配敏感近十倍。
负载也要匹配：同样的长度下，负载电容差 50pF，缓冲延迟可
再差 1--2ns。工程做法因此是三部曲：树先对称，级数对齐，
长度再用蛇形走线补齐。芯片内部的时钟树综合（CTS）做的是
同一件事，只是分支成千上万，由工具在版图上逐级插缓冲、绕
等长，H 形分杈的几何保证从中心到每个角落距离相等。

两种结构并排画出，延迟的累加与抵消一眼可辨。

\begin{pdfdiagram}
  % ===== 左：菊花链 =====
  \node[hw label, font=\footnotesize\bfseries\sffamily] at (2.85,1.6) {菊花链：延迟逐级累加};
  \node[hw control, minimum width=1.0cm] (src) at (0,0) {时钟源};
  \node[hw state, minimum width=0.72cm] (f1) at (1.55,0) {FF};
  \node[hw state, minimum width=0.72cm] (f2) at (2.95,0) {FF};
  \node[hw state, minimum width=0.72cm] (f3) at (4.35,0) {FF};
  \node[hw state, minimum width=0.72cm] (f4) at (5.75,0) {FF};
  \draw[hw bus] (src) -- (f1);
  \draw[hw bus] (f1) -- (f2);
  \draw[hw bus] (f2) -- (f3);
  \draw[hw bus] (f3) -- (f4);
  \node[hw label] at (1.05,0.62) {+6ns};
  \node[hw label] at (2.25,0.62) {+6ns};
  \node[hw label] at (3.65,0.62) {+6ns};
  \node[hw label] at (5.05,0.62) {+6ns};
  \node[hw label] at (0,-0.72) {到达};
  \node[hw label] at (1.55,-0.72) {+6ns};
  \node[hw label] at (2.95,-0.72) {+12ns};
  \node[hw label] at (4.35,-0.72) {+18ns};
  \node[hw label] at (5.75,-0.72) {+24ns};
  \node[hw label, align=center] at (2.85,-1.4) {每级走线加缓冲约 6ns\\8 级链末端比首端晚 48ns};
  % ===== 右：等长树 =====
  \node[hw label, font=\footnotesize\bfseries\sffamily] at (9.35,1.85) {等长树：各分支同时到达};
  \node[hw control, minimum width=1.0cm] (rt) at (7.1,0) {时钟源};
  \node[hw compute, minimum width=0.85cm] (b1) at (8.9,0.75) {BUF};
  \node[hw compute, minimum width=0.85cm] (b2) at (8.9,-0.75) {BUF};
  \node[hw state, minimum width=0.72cm] (g1) at (10.7,1.12) {FF};
  \node[hw state, minimum width=0.72cm] (g2) at (10.7,0.38) {FF};
  \node[hw state, minimum width=0.72cm] (g3) at (10.7,-0.38) {FF};
  \node[hw state, minimum width=0.72cm] (g4) at (10.7,-1.12) {FF};
  \draw[hw bus] (rt.east) -- (8.1,0) |- (b1.west);
  \draw[hw bus] (8.1,0) |- (b2.west);
  \draw[hw bus] (b1.east) -- (9.75,0.75) |- (g1.west);
  \draw[hw bus] (9.75,0.75) |- (g2.west);
  \draw[hw bus] (b2.east) -- (9.75,-0.75) |- (g3.west);
  \draw[hw bus] (9.75,-0.75) |- (g4.west);
  \node[hw label, align=center] at (9.6,-1.75) {各分支走线等长、级数相同\\所有叶节点同时到达};
\end{pdfdiagram}
\diagramnote{时钟分配的两种结构。菊花链中每级走线与缓冲延迟逐级累加，末端沿越来越晚；等长树让每条"根到叶"路径的走线长度与缓冲级数相同，缓冲延迟在分支间相消，所有触发器在结构上同时看到时钟沿。}

树形还有第二个理由：驱动能力。一个时钟源要送到几十上百个
触发器，负载电容累加起来是几十到几百 pF，单级门驱动这样
的负载，延迟从 5ns 涨到几十 ns，边沿也变缓——缓边沿本身
又制造新的偏斜与保持问题。逐级缓冲让每个缓冲器只面对本级
扇出的小负载，延迟可预测、边沿保持陡峭。板级实现里，专用
时钟缓冲芯片（一路进、多路等相位出）就是一棵预制好的小树。

树形分配的代价是功耗。时钟网络是全芯片活动率最高的网络：
每个周期无条件全翻转，负载电容又是各网络中最大的一类，现
代处理器里时钟树可占动态功耗的三到四成。这正是低功耗设计
给时钟树配门控的原因——而门控必须由时钟树工具统一插入并
重新平衡，手工在某一枝上加门，偏斜立刻失控。功耗与偏斜在
这里短兵相接：省功耗的手段制造偏斜，治偏斜的手段费功耗，
工程上是同一个权衡的两面。

菊花链并非一无是处：移位寄存器、扫描链这类数据本身就逐级
流动的结构，时钟顺着数据方向走反而有利——下一级时钟更晚，
恰好帮助保持约束。但它服务的是"顺序"而非"同时"；凡要求
沿对齐的地方，答案都是树。

测量偏斜用双通道示波器：两路探头分别搭在两个触发器的时钟
脚，直接读沿的时间差。注意探头必须同型号、同长度、同衰减
比——探头差异本身会冒充偏斜。低速教学机上 1ns 以下的偏斜，
很容易被测量装置自己的不对称掩盖，这也是一条测量教训。

偏斜可以在面包板上亲手制造。搭一个两位移位寄存器：第一级
的时钟直接取自时钟源，第二级的时钟故意绕一条远路，或串两
级反相器充当延迟。缓慢送入已知图样，观察第二级输出：第二
级时钟晚到得足够多时，同一个时钟沿会被两级先后采到，数据
连跳两位——保持违例的实物版。把绕路拆掉、两级等长，故障
消失。这个几分钟的实验比任何公式都更能让人记住："时钟沿
必须同时到达"不是修辞，是可以被违反、也必然会出事的物理
约定。

\begin{longtable}{L{0.2\linewidth}L{0.32\linewidth}L{0.36\linewidth}}
\toprule
关键问题 & 菊花链 & 等长树（逐级缓冲） \\
\midrule
沿到达时刻 & 逐级递晚，末端晚 N 段延迟 & 结构上同时到达 \\
偏斜量级 & 8 级链约 48ns，超周期 & 长度差 10cm 约 0.67ns \\
驱动能力 & 首级驱动全部负载 & 每级缓冲只驱动本级扇出 \\
适用结构 & 移位、扫描等顺序型数据流 & 一切要求沿对齐的同步系统 \\
\bottomrule
\end{longtable}

\topic{偏斜吃掉时序预算}第二章的预算例：周期 20ns（50MHz），
\code{tCQ_max} = 3ns，组合逻辑最坏 9ns，\code{tSU} = 2ns，
无偏斜时占用 14ns，余量 6ns。现在把偏斜按方向分两次代入，
结论会分裂成两条。

建立检查怕数据来得晚。若目标触发器的时钟比源早到 1ns，有
效周期只剩 19ns：14ns 需求对 19ns 供给，余量降到 5ns。同
一偏斜换个方向——目标时钟晚到 1ns——有效周期变成 21ns，
余量升到 7ns。对建立而言，"目的钟晚到"是礼物。

保持检查怕数据走得早。目标触发器采样沿之后，旧数据必须再
坚持 \code{tH} = 0.5ns。若目标时钟晚到 1ns，坚持时间被延
长成 $0.5 + 1 = 1.5$ns；而数据最早在 \code{tCQ_min} = 1.5ns
后就可能更新——1.5 对 1.5，余量为零，站在违例边缘。偏斜
到 2ns，要求 2.5ns 对供给 1.5ns，保持违例 1ns——这时把周
期拉长到 100ns 也无济于事，因为保持约束与周期无关。刚才对
建立的"礼物"，在保持上是债务：同一偏斜方向相反时，两种
后果互换。这就是"偏斜吃掉预算"的准确含义：它不只是少了
几 ns，而是同时从两个方向威胁两条约束。

\begin{longtable}{L{0.14\linewidth}L{0.24\linewidth}L{0.32\linewidth}L{0.18\linewidth}}
\toprule
检查 & 不利偏斜方向 & 代入数字 & 结果 \\
\midrule
建立 & 目标时钟早到 1ns & 有效周期 19ns 对需求 14ns & 余量 5ns \\
保持 & 目标时钟晚到 1ns & 供给 1.5ns 对要求 1.5ns & 余量 0，边缘 \\
保持 & 目标时钟晚到 2ns & 供给 1.5ns 对要求 2.5ns & 违例 1ns \\
\bottomrule
\end{longtable}

修复手段也分两条。建立问题可以降频换时间，因为建立余量随
周期增长；保持问题与频率无关，只能在数据路径补最小延迟
（插缓冲器、绕一段线）或修时钟树本身。高级设计里还有"有
用偏斜"（useful skew）：故意让关键路径的目的钟晚到一点，
向保持还富裕的路径借时间——这是把双刃剑，借来的每一分都
要在保持端偿还。排查看错误指纹：随频率改变的错是建立方向，
不随频率改变的错先查保持。

偏斜还有一个隐蔽来源：时钟门控。用与门把时钟"关掉"以省
电或实现使能，门本身的几 ns 延迟只出现在被门控的分支上，
等于在树的一侧凭空加了一级缓冲——偏斜立刻超标。第二章的
结论在这里仍然成立：教学机上不要门控时钟，用时钟使能
（CE）让触发器自持；门控时钟是带完整时钟树工具链的设计里
才允许的手法。

把第二章那张预算表补全：占用项依次是 \code{tCQ_max} 3ns、
组合 9ns、\code{tSU} 2ns、偏斜 1ns、抖动窗口 0.7ns，合计
15.7ns；周期 20ns，余量 4.3ns。偏斜与抖动单独看都不致命，
写进同一行才知道预算还剩多少——这就是时序预算必须把所有
项列全的原因。

\topic{抖动是沿位置的随机游移}偏斜是结构的账，结构定了它
就不变；抖动是噪声的账，每个沿都在理想位置附近游移。按观
察窗口分三类：\hwtiming{周期抖动}，单个周期长度对理想周
期的偏离；\hwtiming{周期间抖动}，相邻两个周期长度之差；
\hwtiming{长期抖动}，沿位置对理想位置的累积偏离（TIE）。
来源在教学机上就能找齐三个：电源噪声——门延迟对电压敏感，
量级上 1\% 的电压变化带来 1\% 级的延迟变化，100mV 电源波
动（5V 的 2\%）扫过 5ns 的缓冲链，沿位置就被推着走 100ps
量级；串扰——邻近信号翻转通过容性、感性耦合在时钟线上叠
加台阶，§四展开；振荡器与缓冲链自身的热噪声与闪烁噪声，
这是再干净的板子也去不掉的底。

随机抖动的统计描述用高斯模型：沿位置以理想位置为中心，标
准差 $\sigma$ 即均方根（RMS）抖动。"峰值抖动"对高斯分布
没有绝对上限，只能按可容忍的失误率取倍数：$\pm$3$\sigma$
覆盖 99.73\%，$\pm$7$\sigma$ 才把失误率压到 $10^{-12}$ 量
级。一只 RMS 50ps 的时钟按 $\pm$7$\sigma$ 取窗口，沿位置
的不确定范围是 $\pm$350ps，峰峰 0.7ns。写进预算的做法很
直接：建立检查用最坏周期——标称周期减去抖动窗口。与上一
节的 1ns 偏斜合并：20ns 周期、14ns 需求，余量从纸面的 6ns
被偏斜和抖动合计吃掉 1.7ns，剩 4.3ns。确定性抖动（占空比
失真、与数据码型相关的分量）不是高斯的，按峰峰值与随机分
量相加：总抖动 = DJ + $n \times$ RJ。多级缓冲链上，各级随
机抖动按平方和开根累积，所以时钟链每多一级，抖动底就厚一
分。

三种抖动各有主管的场合。周期抖动直接决定最坏周期，是同步
时序预算的常客；周期间抖动刻画相邻周期的突变，PLL 与延迟
锁定环（DLL）能否跟得上，看它；长期抖动决定串行链路和多
周期同步的取样窗口，通信系统先问它。预算时选错种类等于白
算：给 50MHz 同步机用长期抖动指标，或给串行链路只验周期
抖动，都是把噪声看漏了一半。

测量抖动最直观的是示波器余辉（persistence）模式：让波形
持续叠加几十秒，时钟沿"涂抹"出的宽度就是抖动的峰峰直观
值；要求定量时，用时间间隔误差（TIE）直方图读出 RMS。抖
动与 §一的电源噪声是同一枚硬币的两面，测量时可以直接验证：
两路探头一路看电源轨纹波、一路看时钟沿的散开，翻转密集、
纹波增大时沿的散开跟着变宽——电源噪声调制门延迟的链条就
此闭环。反过来，给时钟缓冲器就近补一只 100nF，沿的散开会
明显收窄。这是少数能用一块面包板直接演示的"噪声变抖动"
实验。

\begin{longtable}{L{0.16\linewidth}L{0.28\linewidth}L{0.24\linewidth}L{0.2\linewidth}}
\toprule
类型 & 定义 & 主要来源 & 预算写法 \\
\midrule
周期抖动 & 单周期长度对理想值的偏离 & 电源噪声、缓冲链噪声 & 用最坏周期代入建立检查 \\
周期间抖动 & 相邻两周期长度之差 & 电源尖峰、串扰事件 & 限制 PLL 跟踪假设 \\
长期抖动（TIE） & 沿位置对理想位置的累积偏离 & 振荡器漂移、低频噪声 & 与偏斜并列扣除余量 \\
\bottomrule
\end{longtable}

把同一个时钟沿用余辉模式叠加几十秒，沿位置的游移就涂抹成一条可见的带，带的宽度与统计指标直接对应。

\begin{waveform}[x=0.9cm,y=0.9cm]
  % 余辉带：实际沿在理想位置两侧游移
  \fill[BookAccent!14] (3.55,1.0) rectangle (4.45,2.6);
  \draw[BookAccent, line width=0.9pt] (0.6,1.0) -- (3.55,1.0) -- (3.55,2.6) -- (6.2,2.6);
  \draw[BookAccent, line width=0.9pt] (1.4,1.0) -- (4.45,1.0) -- (4.45,2.6) -- (7.6,2.6);
  \draw[BookAccent!75, line width=0.7pt] (3.75,1.0) -- (3.75,2.6);
  \draw[BookAccent!75, line width=0.7pt] (4.25,1.0) -- (4.25,2.6);
  % 理想沿
  \draw[BookMuted, dashed, line width=0.6pt] (4.0,0.72) -- (4.0,3.05);
  \node[hw label] at (4.0,3.28) {理想沿位置};
  % 峰峰抖动
  \draw[{Stealth[length=1.8mm]}-{Stealth[length=1.8mm]}, BookRed, line width=0.6pt] (3.55,0.55) -- (4.45,0.55);
  \node[hw label] at (4.0,0.24) {峰峰抖动（按 $\pm 7\sigma$ 取窗口）};
  % RMS 抖动
  \draw[{Stealth[length=1.8mm]}-{Stealth[length=1.8mm]}, BookRed, line width=0.6pt] (3.82,2.88) -- (4.18,2.88);
  \node[hw label, anchor=west] at (4.55,2.88) {$\pm\sigma$（RMS 抖动）};
  \node[hw label, anchor=east] at (0.4,1.8) {时钟沿（余辉叠加）};
\end{waveform}
\diagramnote{抖动的直接图像：示波器余辉模式下，同一个时钟沿在理想位置两侧涂抹成一条带。带的散布宽度对应 RMS 抖动 $\sigma$，时序预算按 $\pm 7\sigma$ 的峰峰窗口扣除。}

\topic{锁相环：倍频、去偏斜和滤波}PLL 的基本结构五环相扣：
鉴相器比较参考沿与反馈沿的先后，输出与相位差成正比的误差
信号；电荷泵把误差换成电流脉冲；环路滤波器（RC 低通）把脉
冲积成平缓的控制电压；压控振荡器（VCO）按控制电压改变频
率；分频器（$\div N$）把输出降回参考频率附近，送回鉴相器。
环路锁定时反馈沿与参考沿对齐，$f_{out} = N \times f_{ref}$。

用途一，倍频。片外 12MHz 晶体，片内乘 8 得 96MHz。为什么
不直接买 96MHz 振荡器？因为 96MHz 方波在 PCB 上的分配本身
就是偏斜、反射与辐射的麻烦（§三、§四）；把高频限制在芯
片内部、板级只走 12MHz，是高速系统的普遍取舍。用途二，去
偏斜。把时钟缓冲器放进反馈环：鉴相器比较的不是 VCO 输出，
而是缓冲器输出，环路锁定强迫"缓冲后的沿"对齐参考沿，缓
冲器那几 ns 延迟被相位比较自动吸收——这叫零延迟缓冲。板
级多片芯片要共用严格对齐的时钟时，专用 zero-delay buffer
芯片就是这个原理。用途三，滤波。环路带宽（例如 1MHz）以下
的参考抖动，环路跟随；带宽以上的快抖动，环路来不及反应，
输出由 VCO 自身的短期稳定度决定——参考上的 10MHz 抖动分
量就这样被滤掉。给抖动的参考"洗干净"，靠的就是把带宽设
低；代价是锁定变慢，上电后要等几十到几百 $\mu$s，复位电
路必须等 PLL 的 locked 信号才能释放。

倍频不是没有代价：相位噪声随倍频比恶化。频率乘 N 的同时，
相位扰动也乘 N，相位噪声谱整体上移 $20\log_{10}N$ dB——
乘 8 就是约 18dB。参考晶体在偏离载波 1kHz 处的相位噪声若
是 -120dBc/Hz，倍频到 96MHz 后同一偏移处约 -102dBc/Hz。
高倍频系统对参考质量的要求因此不降反升：PLL 放大频率的同
时，也放大了参考的相位起伏。这与"调整器"的定位一致——
它传递并放大参考的特性，而不是发明新的特性。

微控制器把这套结构做成了寄存器配置：外部 8MHz 晶体经 PLL
乘 9 得 72MHz 一类的选项，是启动代码里最平常的几行。配置
次序有讲究：先起振、等稳定、再切时钟源、最后等 PLL 锁定
——任何一步抢跑，CPU 就跑在一个还没站稳的时钟上。工业设
计里还会给 PLL 配失锁检测：锁定丢失时自动切回安全时钟，
免得一个时钟毛刺把整机带进未知状态。

定位要准确：PLL 不是时钟发生器，而是时钟调整器。振荡本身
由 VCO 完成，VCO 自由运行时频率随电压、温度漂移；PLL 做的
是把这只不听话的振荡器锁到外部参考上，让输出继承参考的准
确度。参考才是准确度的来源，环路只是调整机构。用 PLL 的第
一问永远是：参考从哪里来，它本身有多准。

\topic{晶振与振荡器的选择}参考源的常规选项有两个。晶体
（quartz crystal）是一块石英薄片，靠压电效应在机械谐振频
率上振动，本身不能起振，需要配振荡电路：教科书结构是皮尔
斯（Pierce）——一只无缓冲反相器（74HCU04，或 MCU 片内的
同等电路）跨接晶体，晶体两端各经一只负载电容到地，再并一
只 M$\Omega$ 级反馈电阻把反相器偏置在线性区。独立振荡器
模块（有源晶振）把晶体、放大与缓冲封在一个四脚金属壳里：
VCC、GND、输出，通电即出方波，即插即用。

晶体的坑集中在负载电容。数据手册给的是负载电容规格 CL，
常见 18pF，它等于两只外接电容的串联值加走线杂散：
$CL = C_1 C_2 / (C_1 + C_2) + C_{stray}$。取 $C_1 = C_2 =
27$pF、杂散按 5pF 估：$27 \times 27 / 54 = 13.5$pF，加 5pF
得 18.5pF，与规格吻合。配不准不会停振，但频率被拉动几十
ppm：32.768kHz 钟表晶体偏 20ppm，一天走差 $86400 \times
20 \times 10^{-6} \approx 1.7$s。起振也有讲究：晶体 Q 值高，
幅度要 ms 级才能建立，上电后最初几 ms 时钟不可用，复位电
路必须等它。布局上晶体要紧贴芯片、下方不走线——晶体走线
是全板最敏感的高阻节点，串扰直接变成频率抖动。

两类常见晶体的差异也值得知道：MHz 级晶体多为 AT 切石英片，
等效电阻几十 $\Omega$，起振快；32.768kHz 钟表晶体是音叉结
构，等效电阻几十 k$\Omega$，驱动电平必须极低，过驱反而缩
短寿命甚至停振。给 RTC 电路配晶体时，串联一只限流电阻、
负载电容严格按规格，是"能振"与"振得准"的分界线。

数据手册里三个相近的词要分开：准确度指出厂时（25 $^{\circ}$C、
额定负载下）频率与标称值的偏差；稳定度指温度、电压变化范
围内频率的漂移；老化指谐振频率随时间的缓慢单向漂移，典型
每年几 ppm。一只"$\pm$20ppm"的普通晶体通常是三项的总和；
温补晶振（TCXO）用温度补偿电路把稳定度一项收到 $\pm$1ppm
量级，代价是价格与功耗。选型时先问哪一项伤害应用：实时钟
怕老化与温漂，通信接口怕总偏差，按需买项，不必为用不到的
指标付钱。

\begin{longtable}{L{0.2\linewidth}L{0.34\linewidth}L{0.34\linewidth}}
\toprule
关键问题 & 晶体 + 片内振荡电路 & 独立振荡器模块 \\
\midrule
外围与调试 & 需皮尔斯电路与负载电容匹配 & 通电即用，无外围 \\
启动时间 & ms 级，幅度缓慢建立 & ms 级，含内部稳定时间 \\
频率稳定度 & $\pm$20--50ppm，受匹配影响 & 同等级晶体，出厂已校准 \\
成本与引脚 & 器件便宜，占两个引脚 & 单价高，体积大，即插即用 \\
适用场合 & 批量产品、成本敏感设计 & 面包板、调试期、小批量 \\
\bottomrule
\end{longtable}

低速教学机的实用建议：调试期直接用振荡器模块或 555——教
学机不与外部设备对表，ppm 无所谓，555 的 RC 漂移（百分之
几）完全够用。什么时候必须上晶体？当时钟要和外部世界对齐：
UART 串口要求收发两端波特率总误差控制在约 $\pm$2\% 以内，
RC 振荡器已贴近边缘，晶体 $\pm$20ppm 则富余几个数量级。
判据只有一句：只跟自己走的机器，频率差不多就行；要跟别的
设备对表的机器，频率必须是有约定的 ppm。

振荡器模块也不是插上就万事大吉：它自己是耗电与时钟驱动的
大户，电源脚同样需要 100nF 去耦；带 OE 脚的模块注意使能电
平，悬空可能停在半导通状态；输出驱动长线时同样受 §三反射
问题约束，必要时串一只几十 $\Omega$ 的源端电阻。时钟源越
"成品化"，越容易让人忘记它仍是一个模拟与数字混合的电路。

排查晶体问题的顺序：不起振，先查负载电容是否配错（过大难
振、过小频偏），再查是否误用了带缓冲的反相器——74HC04
不行，要 74HCU04 这类无缓冲型，缓冲级的额外相移会破坏起
振条件——最后查晶体本身是否受摔受损。频率用频率计或示波
器直读；频率对而系统偶发错，查振幅与驱动电平是否不足。

\begin{keyidea}
时钟分配的核心是管理两类时间差：偏斜是确定的、由结构决定，
用对称的树把它收小；抖动是随机的、由噪声决定，用统计窗口
把它计入预算。锁相环与晶振管的是源头质量：一个负责把振荡
调整到与参考对齐，一个负责提供值得对齐的参考。
\end{keyidea}

\section{三、信号完整性：传输线、反射与端接}

本节的出发点只有一句话：边沿不是瞬间到达的。一个边
沿从驱动器出发，要以有限速度沿走线推进，途中每遇到
一次阻抗变化，就分出一部分波反射回去。本节按问题出
现的次序组织：先定“多长的走线算长”，再看长线上的
行波由什么参数刻画，然后演算反射怎样形成振铃，最后
给出端接的三种标准做法，以及过冲、回沟对输入端的压
力与处置。全节只用三个量：传播速度、特性阻抗、反射
系数；每一条结论都能还原成这三者的四则运算。

\topic{短线判据的定量版}电信号在导线里推进的载体是
电磁场，速度由介质决定：$v=c/\sqrt{\varepsilon_r}$。
真空里 $c\approx30$cm/ns；FR4 板材的相对介电常数约
4.4，内层带状线的速度就是 $30/\sqrt{4.4}\approx14.3$cm/ns；
表层微带线一半贴着空气，有效介电常数降到 3 上下，
速度约 16--17cm/ns。工程估算统一取 15cm/ns，即每厘
米约 67ps：一根 10cm 走线，边沿从一头到另一头要约
0.7ns。这个“慢”是否要紧，取决于和边沿时间的比。

判据比的不是绝对长度，而是线延迟与边沿时间的比值。
边沿沿线推进时，反射波与入射波在途中叠加；若线延迟
远小于边沿时间，反射还没来得及与入射分离，就被淹没
在边沿的斜坡里，线上各点几乎一起动——这是“短”线。
经验分界线取单向延迟等于边沿时间的六分之一：超过它，
反射波在边沿尚未走完时已独立成形，过冲与振铃开始可
以分辨，必须按传输线处理。写成算式：临界长度
$l_c=t_r\cdot v/6$，$t_r$ 为边沿时间。六分之一不是
物理定律，是“反射刚开始能被看见”的工程刻度，取四
分之一或八分之一的主张都有，结论差不到一倍。

代入数字。74HC 在 5V 下边沿约 5--8ns，
$l_c=5\text{ns}\times15\text{cm/ns}/6\approx12.5$cm，
慢边沿对应 20cm——面包板上十几厘米的飞线大多仍在
短线区，这正是前面七章能用集总模型讲下来的资格。换
成边沿 1ns 的高速 CMOS，临界长度只剩 2.5cm；边沿
0.3ns 的现代 FPGA 输出，5mm 就算长线。同一块板子上，
“这根线短不短”没有统一答案，要一根一根拿器件的边
沿去比。

短线区与长线区是两套模型。短线区用集总模型：整条线
当作一个节点，寄生参数并成一个对地电容（10cm 走线
约 5--15pF），驱动器只是多带了一个负载——前面各章
默认的就是这个世界。长线区必须用分布模型：线上每一
点的电压电流可以不同，要用单位长度的电感 $L'$ 与电
容 $C'$ 刻画，“波”的概念登场。

\begin{longtable}{L{0.2\linewidth}L{0.15\linewidth}L{0.22\linewidth}L{0.29\linewidth}}
\toprule
逻辑系列 & 典型边沿 & 临界长度 $t_r v/6$ & 10cm 走线怎么算 \\
\midrule
74HC（5V） & 5--8ns & 12--20cm & 短线，集总模型足够 \\
74AC、74AHC & 2--3ns & 5--7.5cm & 临界附近，按长线稳妥 \\
现代 FPGA、处理器 IO & 0.1--1ns & 2.5mm--2.5cm & 长线，必须按传输线设计 \\
\bottomrule
\end{longtable}

不知道边沿多快时有两条路。查数据手册的上升/下降时
间一项（74HC 手册给的是 50pF 负载下的典型值）；或
用示波器直接量，注意示波器加探头的等效带宽要满足
$0.35/t_r$——量 2ns 的边沿，测量链至少要有
175MHz，即 200MHz 这一档。既没有手册也没有示波器
时，按“线长接近临界值就当长线”设计：为长线做的端
接用在短线上无害，反过来则直接是故障，宁保守勿侥幸。

\topic{特性阻抗不是电阻}边沿刚进入走线时，并不知道
末端接了什么——它只看到沿途每一小段要充的电和要建
立的磁场。设线每厘米电容 $C'$、传播速度 $v$，波前
每推进 1cm 用时 $1/v$，并把这一厘米的电容充到 $V$，
需电荷 $Q=C'V$，于是波前从驱动器抽走的电流
$I=Q/t=C'vV$。电压与电流之比 $V/I=1/(C'v)$ 只由线
本身决定，定义为特性阻抗 $Z_0$；用 $v=1/\sqrt{L'C'}$
代入，就得到常见的形式 $Z_0=\sqrt{L'/C'}$。它的物
理含义一句话说尽：行波看到的瞬时负载。波在途中的一
切行为——分压、反射、吸收——都由这个比值决定，与
末端还有多远无关。

代入典型值。一条 50$\Omega$ 微带线，$C'\approx1.3$pF/cm，
$L'\approx3.3$nH/cm。验算：
$Z_0=\sqrt{3.3\text{nH}/1.3\text{pF}}\approx\sqrt{2538}\approx50\Omega$；
$v=1/\sqrt{L'C'}=1/\sqrt{3.3\text{nH}\times1.3\text{pF}}\approx65$ps/cm，
即约 15.3cm/ns，与上一段的取值自洽。3.3V 边沿进入
这样的线，波前瞬时电流 $I=3.3/50\approx66$mA——哪
怕末端开路、稳态电流为零，驱动器在边沿期间也要供出
这 66mA。这就是“传输线负载”与“电阻负载”最直观
的差别：稳态开路，瞬态却是 50$\Omega$。

50$\Omega$ 不是物理常数，而是历史折中的沉淀。同轴
电缆时代，空气介质同轴衰减最小的点在约 77$\Omega$，
功率容量最大的点在约 30$\Omega$，50$\Omega$ 是两者
的折中，随射频与仪器工业成为标准；视频与有线电视继
承了低损耗那一端，75$\Omega$ 沿用至今。PCB 沿用
50$\Omega$ 还有现实理由：现代 CMOS 驱动器输出阻抗
十欧上下，串一只三十几欧的电阻就能配到 50$\Omega$；
定得更高驱动器够不着，定得更低静态功耗又太大。双绞
线约 100$\Omega$、高速串行链路常用 85--100$\Omega$
差分，都是同一类折中的不同选择。

三个“无关”值得逐一确认。与线长无关：$L'$、$C'$
都是每单位长度的参数，比值里长度被约掉。与频率无关：
无损近似下两者都不含频率，$Z_0$ 在很宽频段里是常数
——它不是电阻，却表现得像一个电阻，这正是能拿电阻
做端接的根据。与绝对粗细无关：把横截面按同一比例放
大，$L'$ 与 $C'$ 的比不变，$Z_0$ 不变——它只由几
何比例（线宽对介质厚度之比）与介电常数决定。但若只
加宽线、不动介质，$C'$ 增、$L'$ 减，$Z_0$ 按平方根
下降：50$\Omega$ 的线线宽加倍，阻抗大致减半。电源
线加宽总是好事，信号线加宽之前要先算这笔数。

\begin{warningbox}
用万用表电阻档量不出 50$\Omega$：欧姆表量的是导体
的直流电阻，一条铜走线只有毫欧级。特性阻抗只对变化
中的边沿“显形”，稳态下它不存在。反过来也别走到
另一极端，把传输线想成神秘元件：它就是一段处处一样
的介质，全部学问只有 $L'$、$C'$ 两个参数和它们的比。
\end{warningbox}

测量上，没有仪器也能估 $Z_0$：用一只电容表量一段已
知长度走线对参考面的总电容，除以长度得 $C'$，再由
$Z_0=1/(C'v)$、取 $v\approx15$cm/ns 反推。例：30cm
走线量得 40pF，$C'=1.33$pF/cm，
$Z_0=1/(1.33\text{pF/cm}\times15\text{cm/ns})\approx50\Omega$。
专业仪器是时域反射计（TDR）：向走线发一个已知快边
沿，从回波形状直接读出沿线每处的阻抗剖面，连接器、
过孔、线宽突变在图上各成一个台阶——原理就是本节三
个公式的反向使用。

\topic{反射系数与振铃演算}行波到达阻抗突变处，波带
来的电压电流比是 $Z_0$，新阻抗要求的却是另一个比
值，矛盾只能靠再产生一个波调和：一部分透射，一部分
反射。由边界上电压连续、电流连续两个条件，解出反射
电压与入射电压之比——反射系数
$\Gamma=(Z_L-Z_0)/(Z_L+Z_0)$。三个特例值得记住：
$Z_L=Z_0$ 时 $\Gamma=0$，波被完整吸收，这叫匹配；
末端开路 $Z_L\to\infty$ 时 $\Gamma=+1$，反射与入射
同号，末端电压跳到入射波的两倍；末端短接 $Z_L=0$
时 $\Gamma=-1$，反射反号，末端电压归零、电流加倍。

反射波回到源端，源端同样是一处阻抗突变，按
$\Gamma_S=(Z_S-Z_0)/(Z_S+Z_0)$ 再反射一次。CMOS
驱动器的输出阻抗通常小于 $Z_0$，$\Gamma_S$ 为负：
回波到源端被反号弹回，再向末端走去，到末端又全反
射……波就在两端之间来回弹，每弹一圈，幅度打一个
$|\Gamma_S\Gamma_L|$ 的折扣。振铃不是神秘现象，它
就是这场来回弹的直接记录。

完整演算一遍。3.3V 系统、50$\Omega$ 走线、源阻
10$\Omega$、末端开路，记单向延迟为 $T$。入射波幅度
由源阻与线阻抗分压：$3.3\times50/(50+10)=2.75$V。
$t=T$ 时波到末端，$\Gamma_L=+1$，末端电压从 0 跳到
$2.75+2.75=5.5$V——稳态 3.3V 的 1.67 倍；若源阻为
零，入射即满幅 3.3V，末端将冲到 6.6V，接近电源电
压的两倍，这就是“开路全反射，过冲接近两倍”的准确
含义。回波 2.75V 返到源端，
$\Gamma_S=(10-50)/(10+50)=-2/3$，弹回 $-1.83$V；
$t=3T$ 时它到达末端并再次全反射，末端电压
$5.5-1.83-1.83=1.83$V。如此往复，每往返一圈幅度乘
2/3，逐拍收敛到 3.3V：

\begin{longtable}{L{0.1\linewidth}L{0.36\linewidth}L{0.14\linewidth}L{0.24\linewidth}}
\toprule
时刻 & 到达末端的波 & 末端电压 & 现象 \\
\midrule
$t=T$ & 入射 2.75V，末端全反射再加 2.75V & 5.5V & 最大过冲 \\
$t=3T$ & 源端弹回 $-1.83$V，末端再反射 & 1.83V & 最深回沟 \\
$t=5T$ & $+1.22$V 到达并翻倍 & 4.28V & 二次过冲 \\
$t=7T$ & $-0.81$V 到达并翻倍 & 2.65V & 二次回沟 \\
$t=9T$ & $+0.54$V 到达并翻倍 & 3.73V & 向 3.3V 收敛 \\
\bottomrule
\end{longtable}

两点解读。其一，末端电压每 $2T$ 跳一次，一个完整的
“峰—谷—峰”历时 $4T$：15cm 走线 $T\approx1$ns，
振铃约 250MHz——线越短振铃越高频，低速示波器看不
见，看不见不等于不存在。其二，每圈衰减因子
$|\Gamma_S\Gamma_L|=2/3$ 由两端不匹配程度的乘积决
定；源阻若只有 1$\Omega$（强驱动器），$\Gamma_S$
约为 $-0.96$，振铃要拖很多圈才平息。把这套逐拍递推
画成“时间—位置”斜线图（弹跳图），就是分析任意端
接组合的通用工具：每个交点按当地的 $\Gamma$ 分裂，
照抄上面的四则运算即可。

把这张逐拍表画成末端电压波形，振铃的形状与收敛过程直接可见。

\begin{waveform}[x=0.86cm,y=0.9cm]
  % 稳态参考线 3.3V
  \draw[hw return] (0,1.82) -- (10.4,1.82);
  \node[hw label, anchor=west] at (10.5,1.82) {稳态 3.3V};
  % 跳变时刻参考虚线
  \draw[BookMuted, dashed, line width=0.5pt] (1,0) -- (1,3.3);
  \draw[BookMuted, dashed, line width=0.5pt] (3,0) -- (3,3.3);
  \draw[BookMuted, dashed, line width=0.5pt] (5,0) -- (5,3.3);
  \draw[BookMuted, dashed, line width=0.5pt] (7,0) -- (7,3.3);
  \draw[BookMuted, dashed, line width=0.5pt] (9,0) -- (9,3.3);
  \node[hw label, anchor=north] at (1,-0.08) {$T$};
  \node[hw label, anchor=north] at (3,-0.08) {$3T$};
  \node[hw label, anchor=north] at (5,-0.08) {$5T$};
  \node[hw label, anchor=north] at (7,-0.08) {$7T$};
  \node[hw label, anchor=north] at (9,-0.08) {$9T$};
  % 时间轴
  \draw[line width=0.55pt, -{Stealth}] (0,0) -- (10.4,0);
  % 末端电压波形（3.3V 系统、50Ω 走线、源阻 10Ω、末端开路）
  \draw[BookAccent, line width=0.95pt] (0,0) -- (1,0) -- (1,3.02) -- (3,3.02) -- (3,1.01) -- (5,1.01) -- (5,2.35) -- (7,2.35) -- (7,1.46) -- (9,1.46) -- (9,2.05) -- (10.4,2.05);
  \node[hw label, anchor=east] at (-0.2,1.5) {末端电压};
  \node[hw label, anchor=south] at (2,3.08) {5.5V 过冲};
  \node[hw label, anchor=north] at (4,0.95) {1.83V 回沟};
  \node[hw label, anchor=south] at (6,2.4) {4.28V};
  \node[hw label, anchor=north] at (8,1.38) {2.65V};
  \node[hw label, anchor=south] at (9.7,2.1) {3.73V};
\end{waveform}
\diagramnote{3.3V、50$\Omega$ 走线、源阻 10$\Omega$、末端开路时的末端电压逐拍演算。入射 2.75V 在 $t=T$ 到达末端并全反射成 5.5V，回波在源端被负反射后逐拍抵消，每往返一圈幅度乘 $2/3$，最终收敛到稳态 3.3V。}

测量上反过来用这套表：在末端量出振铃周期 $4T$，就
能反推走线的单向延迟与等效长度——这是没有 TDR 时
的土法测长。若末端改短接，同一工具给出源端波形：先
是 2.75V 的台阶，$2T$ 后回落，逐拍趋向 0；稳态就是
源阻对地短路，10$\Omega$ 源阻下电流 330mA——这也
是拿短接输出脚来“试波形”为什么危险的原因。

\topic{端接的三种做法}端接的总原则只有一句：要么让
反射不发生，要么让它只发生一次并被吸收。前者在末端
做匹配，令 $\Gamma_L=0$；后者在源端做匹配，令
$\Gamma_S=0$，故意放一次全反射过去再把它接住。三条
标准路线对应不同的拓扑与功耗预算。

源端串联端接：在驱动器出口串一只电阻，使“驱动器内
阻加串联电阻”等于 $Z_0$——内阻 15$\Omega$ 就串
39$\Omega$。入射波因此只有半幅，沿线走到开路末端全
反射翻倍，恰好一次到达满幅；回波回到源端时被匹配吸
收，$\Gamma_S\approx0$，故事就此结束，稳态功耗为零。
代价藏在波形里：线的中途各点看到的是“先半幅、回波
经过再到满幅”的两段台阶，只有末端那一点的边沿干净。
所以它只适合点对点、负载集中在末端——时钟线、片选
线的首选；负载沿线分布时，中途的接收器会在阈值附近
看到一个停顿平台，这不能接受。顺带一记：74HC 输出
级导通电阻典型几十欧，本身就近似一只天然的串联端接
电阻，面包板上 74HC 的振铃常比理论值温和——这是运
气，不是可以依赖的设计。

末端并联端接：在线的最远端接一只 $R_L=Z_0$ 的电阻
到地，$\Gamma_L=0$，入射波一次到达即满幅，全程无反
射，沿线每一点看到同一个干净边沿——多负载沿线分布
的总线只能走这条路。代价是功耗：3.3V 驱动 50$\Omega$，
高电平期间每线恒流 66mA，八位总线全高就是半安培级，
驱动器与电源都要按这个电流设计。这是拿功耗换干净的
典型交易。

戴维南端接用两只电阻：上拉、下拉各取 $2Z_0$（
50$\Omega$ 线各用 100$\Omega$），并联等效恰好
50$\Omega$，中点电压 1.65V 顺手给总线一个确定的空
闲电平——需要空闲态有定义的总线用它，静态功耗比单
电阻小些但仍可观。AC 端接把端接电阻串一只
100--200pF 电容再到地：直流被隔断，静态功耗近零；
边沿到来的几纳秒内电容近似短路，电路呈现 $Z_0$。代
价是每次翻转要给这只电容充放电，且效果对边沿快慢略
有依赖——它适合连续翻转的时钟，不适合偶尔跳一下的
控制线。

\begin{longtable}{L{0.13\linewidth}L{0.23\linewidth}L{0.24\linewidth}L{0.13\linewidth}L{0.15\linewidth}}
\toprule
端接方式 & 接法 & 反射行为 & 静态功耗 & 适用场合 \\
\midrule
源端串联 & 出口串 $R=Z_0-Z_S$ & 末端全反射一次后被吸收 & 零 & 点对点，负载在末端 \\
末端并联 & 末端 $R=Z_0$ 到地 & 无反射 & 每线数十 mA & 多负载总线末端 \\
戴维南 & 上拉下拉各 $2Z_0$ & 无反射，空闲电平确定 & 中等 & 需定义空闲态的总线 \\
AC 端接 & $Z_0$ 串 100--200pF 到地 & 边沿期间近似无反射 & 近零 & 连续时钟，功耗敏感 \\
\bottomrule
\end{longtable}

三种做法画成电路对照，电阻的取值依据与等效阻抗如下。

\begin{pdfdiagram}
  % ===== 源端串联 =====
  \node[hw label, font=\footnotesize\bfseries\sffamily] at (1.6,1.5) {源端串联};
  \node[hw compute, minimum width=1.15cm] (dA) at (0,0) {驱动器\\ $R_{out}$};
  \node[hw state, minimum width=1.15cm] (rA) at (3.1,0) {接收端\\ （开路）};
  \draw[line width=0.55pt] (dA.east) -- (1.05,0);
  \draw[line width=0.55pt] (1.05,-0.2) rectangle (1.95,0.2);
  \draw[line width=0.55pt] (1.95,0) -- (rA.west);
  \node[hw label] at (1.5,0.42) {$R_s$};
  \node[hw label, text width=3.4cm, align=center] at (1.6,-1.65) {$R_s = Z_0 - R_{out}$\\末端全反射一次，回波被吸收};
  % ===== 末端并联 =====
  \node[hw label, font=\footnotesize\bfseries\sffamily] at (5.9,1.5) {末端并联};
  \node[hw compute, minimum width=1.15cm] (dB) at (4.5,0) {驱动器};
  \node[hw state, minimum width=1.15cm] (rB) at (6.3,0) {接收端};
  \draw[line width=0.55pt] (dB.east) -- (rB.west);
  \draw[line width=0.55pt] (rB.east) -- (7.15,0) -- (7.15,-0.4);
  \draw[line width=0.55pt] (6.97,-0.4) rectangle (7.33,-1.0);
  \draw[line width=0.55pt] (7.15,-1.0) -- (7.15,-1.2);
  \pic at (7.15,-1.2) {hw gnd};
  \node[hw label, anchor=west] at (7.4,-0.7) {$R_t$};
  \node[hw label, text width=3.4cm, align=center] at (5.7,-1.7) {$R_t = Z_0$ 到地\\$\Gamma_L = 0$，全程无反射};
  % ===== 戴维南 =====
  \node[hw label, font=\footnotesize\bfseries\sffamily] at (10.0,1.5) {戴维南};
  \node[hw compute, minimum width=1.15cm] (dC) at (8.55,0) {驱动器};
  \node[hw state, minimum width=1.15cm] (rC) at (10.35,0) {接收端};
  \draw[line width=0.55pt] (dC.east) -- (rC.west);
  \draw[line width=0.55pt] (rC.east) -- (11.15,0);
  \draw[line width=0.55pt] (11.15,0) -- (11.15,0.4);
  \draw[line width=0.55pt] (10.97,0.4) rectangle (11.33,1.0);
  \draw[line width=0.55pt] (11.15,1.0) -- (11.15,1.15);
  \node[hw label] at (11.15,1.32) {Vcc};
  \draw[line width=0.55pt] (11.15,0) -- (11.15,-0.4);
  \draw[line width=0.55pt] (10.97,-0.4) rectangle (11.33,-1.0);
  \draw[line width=0.55pt] (11.15,-1.0) -- (11.15,-1.2);
  \pic at (11.15,-1.2) {hw gnd};
  \node[hw label, anchor=east] at (10.9,0.7) {$2Z_0$};
  \node[hw label, anchor=east] at (10.9,-0.7) {$2Z_0$};
  \node[hw label, text width=3.4cm, align=center] at (9.55,-1.7) {上拉下拉各 $2Z_0$\\并联等效 $Z_0$，空闲电平确定};
\end{pdfdiagram}
\diagramnote{三种端接的拓扑。源端串联故意放一次全反射过去、再在源端把它吸收，适合点对点；末端并联与戴维南让反射根本不发生，适合多负载总线，代价是静态功耗。}

多负载的布线拓扑与端接绑定。菊花链让一条线依次经过
各负载、在末端一次端接，中途负载只引出短桩（桩本身
要短于临界长度，否则桩也成传输线），是总线的现实选
择；DDR 地址线的 fly-by 结构就是它的工程化版本。星
形从驱动器直接分出 $N$ 条支路，驱动器看到的阻抗变
成 $Z_0/N$，分叉口本身又是一次阻抗突变——支路多于
两三条就很难收拾，只能每条支路各自源端串联、按点对
点逐条处理。端接电阻的摆放也有纪律：源端电阻贴驱动
器出口，末端电阻贴最后一个负载之后，远离位置的端接
等于没接。先按负载分布选拓扑，再按拓扑选端接，顺序
不能反。

\topic{过冲与回沟对输入端的压力}上面演算出的 5.5V
不是纸上数字，它真实出现在 3.3V 系统的输入脚上。每
个 CMOS 输入脚都带两只钳位二极管：一只阳极在脚、阴
极在 Vcc，脚电压超过 Vcc 约 0.5--0.7V 时导通；另一
只反向到地，脚电压低于约 $-0.5$V 时导通。它们的本
职是泄放静电，不是处理日常波形。上升沿的 5.5V 过冲
会打开上管；下降沿的镜像过程把脚电压冲到 $-2.2$V
（入射 $-2.75$V 全反射），打开下管。

保护管导通后的电流由线阻抗限制：
$(5.5-3.3-0.7)/50\approx30$mA，每个边沿一次。手册
通常只给这类电流 $\pm20$mA 上下的定额——偶尔超一
点不会立刻坏，但每个时钟沿都打它一次，是按年计的可
靠性消耗。更大的过冲还有急性病：触发芯片内部寄生的
晶闸管结构（闩锁，latch-up），Vcc 与地之间形成低阻
通路，芯片数秒内发烫，只有彻底断电才能解除，反复几
次即永久损伤。闩锁是过冲问题里最不能讨价还价的一种。

回沟的危害在时钟线上最尖锐。数据线上的回沟发生在建
立时间之前很久，采样时早已平息，往往无害；时钟线不
同——接收触发器只认阈值穿越。上面 10$\Omega$ 源阻
的例子里，回沟停在 1.83V，离 1.65V 的典型阈值只剩
约 0.2V 裕量；源阻若只有 5$\Omega$（强驱动），
$\Gamma_S\approx-0.82$，回沟探到约 1.1V，直接跌穿
阈值，随后 $t=5T$ 的回升再向上穿越一次——一次边沿
被数成先下后上两次触发，状态机多走一拍。这类故障表
现为“计数偶尔加二”“程序莫名跳步”，离病根极远，
排查时极易误判为软件问题。

处置按三条思路排：消灭反射（源端串联或末端端接，反
射不存在，回沟无从谈起）；把线变短（放慢边沿或缩短
走线，回到短线判据之内——1MHz 的总线没有理由用
2ns 边沿的器件）；吸收残余（不计速度的信号如复位线，
可在输入端前加 RC：串联 1k$\Omega$、对地 100nF 量级，
把毛刺整体滤掉）。

\begin{longtable}{L{0.15\linewidth}L{0.29\linewidth}L{0.2\linewidth}L{0.22\linewidth}}
\toprule
手段 & 作用机理 & 代价 & 适用场合 \\
\midrule
源端串联 & 末端一次到满幅，无回沟 & 中途点见台阶 & 点对点时钟、片选 \\
末端并联、AC & 反射根本不发生 & 功耗或元件 & 总线、多负载 \\
放慢边沿 & 走线相对边沿重新变短 & 速度上限下降 & 低速总线首选 \\
缩短走线 & 延迟降到临界值以下 & 布局受限 & 一切场合 \\
RC 吸收 & 毛刺被积分滤除 & 边沿变缓 & 复位、低速控制线 \\
\bottomrule
\end{longtable}

测量时记住一件事：传输线上不同位置的波形不同，驱动
脚看到的是入射波与回波的中途叠加，接收脚看到的才是
判决依据——探头必须点在接收脚上。判读两条红线：过
冲顶点不许碰 Vcc$+0.5$V，回沟谷底不许碰阈值区。两
条都碰了还不端接，剩下的就不是会不会坏的问题，而是
什么时候坏的问题。

\section{四、串扰、回流与电磁兼容入门}

§三管的是一根线上的波形；本节管线线之间、芯片内
外、板子内外的相互作用。串扰、回流绕路、同时开关噪
声、差分与辐射，看似五件事，物理上只有一件：变化的
电场与磁场不认走线这个边界。本节按作用范围由近及远
排开，每条规则仍然先给物理图像，再给数量级，最后给
排查方法。

\topic{串扰来自平行长走线}两条平行走线共享同一段介
质，攻击线上的变化有两条路渗进受害线。容性耦合走电
场：攻击线电压变化，经线间互容向受害线注入电流
$I=C_m\,\mathrm{d}V/\mathrm{d}t$，注入电流在受害线
上向两端各分一半。感性耦合走磁场：攻击线电流变化，
互感在受害线上感应电压，驱动起一个与攻击电流反方向
的环流。两种耦合有方向性：朝驱动器方向（近端）去的
两半同相相加，朝末端方向（远端）去的两半反相相消——
近端串扰因此通常最大，带状线里远端甚至接近抵消干净；
微带线因两种耦合的传播速度略有差异，远端会留一点残
余。量串扰，探头先放在受害线的近端。

把两条平行走线的耦合路径画出来，近端与远端毛刺的方向性就清楚了。

\begin{pdfdiagram}
  % 攻击线
  \node[hw compute, minimum width=1.0cm] (drv) at (0,2.4) {驱动器};
  \draw[BookAccent, line width=1.0pt] (drv.east) -- (9.2,2.4);
  \node[hw label, anchor=south west] at (0.6,2.5) {攻击线};
  \draw[-{Stealth[length=2mm]}, BookAccent, line width=0.6pt] (4.6,2.75) -- (6.4,2.75);
  \node[hw label, anchor=south] at (5.5,2.8) {攻击边沿传播};
  % 受害线（含近端、远端毛刺）
  \draw[BookTeal, line width=1.0pt] (0.6,0.9) -- (1.1,0.9) -- (1.1,1.32) -- (1.8,1.32) -- (1.8,0.9) -- (7.7,0.9) -- (7.7,1.14) -- (8.3,1.14) -- (8.3,0.9) -- (9.2,0.9);
  \node[hw label, anchor=north west] at (0.6,0.8) {受害线};
  \node[hw label, anchor=north] at (1.6,0.6) {近端毛刺：同相相加，最大};
  \node[hw label, anchor=north] at (8.0,0.6) {远端毛刺：反相相消，微小};
  % 耦合路径
  \draw[hw return] (3.4,2.4) -- (3.4,0.9);
  \node[hw label, anchor=east] at (3.25,1.65) {互容 $C_m$};
  \draw[hw return] (4.9,2.4) -- (4.9,0.9);
  \node[hw label, anchor=west] at (5.05,1.65) {互感 $L_m$};
  % 3W 间距规则
  \draw[{Stealth[length=1.8mm]}-{Stealth[length=1.8mm]}, BookMuted, line width=0.6pt] (6.6,0.9) -- (6.6,2.4);
  \node[hw label, anchor=west] at (6.75,1.65) {中心距 $\geq 3W$};
\end{pdfdiagram}
\diagramnote{平行走线间的两条耦合路径：互容注入电流、互感感应环流。朝近端去的两半同相相加，近端毛刺最大；朝远端去的两半反相相消，远端只剩残余。中心距拉开到三倍线宽（3W 规则），可把耦合压到两成上下。}

耦合强度随间距下降的速度比直觉快。走线下方有参考面
时，耦合系数近似按 $1/(1+(s/h)^2)$ 衰减，$s$ 为两线
净间距，$h$ 为走线到参考面的介质厚度。取常见情形
$h$ 约等于线宽 $W$：净距为零时耦合记作 100\%，净距
一个线宽约 50\%，净距两个线宽约 20\%。3W 规则——相
邻走线中心距不小于三倍线宽，即净距不小于两倍线宽——
正是把串扰按到两成上下的那条经验线。它不花一个元件，
只花布线面积；继续拉到 5W、10W 收效递减，不如把省
下的面积留给回流。注意 3W 的前提是下方有完整参考面
把场约束住：没有参考面，公式里的 $h$ 失去意义，拉开
多远都不够。

平行长度也有自己的刻度。近端串扰随平行长度增长，直
到平行走线的往返延迟追平边沿时间后饱和，饱和长度
$l_s=t_r\cdot v/2$：6ns 边沿对应约 45cm。面包板上
10cm 的并行走线远没到顶，所以教学机上缩短平行段立
竿见影；反过来，边沿 0.3ns 时饱和长度只剩约 2.3cm，
高速板上几厘米的并行早已到顶——“能容忍多长的并
行”，同样要拿边沿时间去比。

受害线旁边最不能住的三类邻居：时钟、复位、片选（写
使能同理）。时钟边沿快且永不停歇，串扰以时钟周期重
复出现，最容易准时混进数据的采样窗口；复位线上一个
几纳秒的毛刺就可能复位全机，而毛刺窄到普通示波器抓
不住，故障现象（莫名重启）离原因极远；片选或写使能
上的假脉冲意味着假读写——RAM 被改写、外设被误触发，
破坏是永久性的。

\begin{longtable}{L{0.15\linewidth}L{0.42\linewidth}L{0.3\linewidth}}
\toprule
攻击源 & 危险之处 & 布线对策 \\
\midrule
时钟 & 边沿快、周期性重复，易混入采样窗 & 与数据线保持 3W，必要时地线隔离 \\
复位 & 窄毛刺即可复位全机，事后无痕 & 远离快信号、走线短、输入端加 RC \\
片选、写使能 & 假选通等于假读写，数据被破坏 & 与数据线同等管理，注意端接 \\
\bottomrule
\end{longtable}

必须紧邻时，在攻击线与受害线之间布一条接地隔离线，
把电场截到地，耦合明显削弱。但隔离线两端必须就近接
地、长则沿途多点接地；悬空或只接一头的隔离线自己就
是一根转发耦合的天线，比没有更糟。排查串扰有三个可
动的旋钮：毛刺幅度随平行长度增长、随间距加大下降、
随攻击边沿加快而上升——让机器跑起来，把受害线驱动
成固定电平，逐个动这三个变量，符合规律即可确认。若
噪声与“同拍翻转的路数”相关而非与某条邻线相关，那
是后面要讲的 SSN，两者治法不同，先分清再动手。

感性耦合可以用变压器理解：攻击线是初级，受害线是次
级，各一匝，磁路是两者之间共享的那部分磁通——并行
越长、间距越近，共享磁通越多，互感越大；容性耦合则
是一只跨接在两线之间的小电容。两只“寄生元件”的值
都由几何决定，这正是 3W 规则背后那张等效电路。测量
时两种耦合还能分开辨认：近端毛刺的宽度等于并行段的
往返时间，远端毛刺的宽度等于攻击边沿的时间——用这
两个“指纹”，在示波器上就能读出耦合走的是哪条路。

总线场景要把串扰再算重一点。8 位总线并行走线，中间
那根线两侧各有一个攻击者，串扰两份相加——总线里最
惨的总是中间位。最坏的激励是邻位全部同拍同向翻转、
受害位保持不动：两侧的耦合电流同相叠加，等效攻击强度翻倍。
评估一条总线的串扰，要按这个最坏图案算，而不是按平
均翻转率算。

教学机上没有参考面，$h$ 的约束不存在，串扰按最坏情
形估：两根并排 10cm 的杜邦线，互容约 5pF 量级。攻击
线 5V/6ns 的边沿注入电流
$I=C_m\,\mathrm{d}V/\mathrm{d}t\approx5\text{pF}\times5\text{V}/6\text{ns}\approx4$mA，
受害线上两端各分一半：受害线被 50$\Omega$ 驱动时，
2mA 产生的毛刺约 0.1V，尚在噪声容限之内；受害线若
是悬空的输入脚，25pC 的电荷注入几 pF 的节点电容，
毛刺就是伏特级，足以误触发。这就是面包板上“输入脚
不许悬空”的另一半理由：上拉、下拉电阻不仅给逻辑一
个默认值，也是给串扰电流一条低阻旁路。

排线与带状电缆则是把回流影子做成实物：工程排线把地
线按“地—信—信—地”或更密的“地—信—地”间隔排
布，让每根信号旁边都有自己的回流通道，环路面积与串
扰同步受控。一根 26 芯排线里插 5--6 根地线，牺牲的
是几根信号位，换来的是整束线的可靠；全插信号的排线
在高速下几乎一定出错。面包板上把时钟线两侧各插一根
地线，是同一手段的零成本版本。

\topic{回流路径决定环路面积}“电流从地流回去”这句
话在高速下要改写。直流与低频时，回流在整个地平面上
铺开，走电阻最小的路径；高频时回流走电感最小的路径
——恰好是信号走线正下方参考面上的那条窄带，像信号
投在参考面上的影子。物理原因很朴素：贴着信号流，环
路面积最小、环的自感最小，电流总是拣能量最省的路。
推论很重要：高频下的“地”不是一整片等位体，每条信
号线脚下都拖着自己专属的回流影子；影子走不通的地方，
信号一定出问题。

回流被迫绕路有三个经典场景。其一，换层：信号打过孔
换到另一面，若上下参考面之间没有就近的地过孔或层间
通路，回流要绕到远处才能跟上——换层的信号孔旁边配
一只“伴随地孔”，就是为回流修的便桥。其二，跨分割：
参考面被开槽、分区，信号线跨过槽缝，回流被迫绕到槽
的端头，环路面积从“线到面的距离”量级暴涨到“槽长”
量级——地平面分割的代价不在地本身，在于逼着回流画
大圈。其三，连接器与排针：几十根信号的回流若只能挤
过一两根公用地针，公共段就成了所有环路的共享部分，
一根线上的变化经共享段电感耦合进所有线。高速连接器
按信号与地 2:1 到 4:1 插地针，面包板上把地线分散插
满排，都是同一个道理。

环路面积是串扰、辐射、敏感度三者共同的根源：一个环
的电感大体随面积增长（1cm 见方的小环约 10nH 量级），
这个电感既决定它向外辐射多少，也决定它从外界场里捡
起多少、与相邻环互感多少。给共享回流算一笔数：公共
段电感 5nH，8 路信号各 20mA、边沿 5ns 同时经过，
$V=L\,\mathrm{d}I/\mathrm{d}t=5\text{nH}\times160\text{mA}/5\text{ns}=0.16$V
——所有经过这段的信号，低电平被整体垫高 0.16V。环
路面积省下的每一平方毫米，都是三份收益。

查回流路径不需要仪器，只需要一个习惯：对每根关键信
号问一句“它的回流从哪条路回来”，再顺着答案走一遍。
答不上来的、要绕远的、要跟别人挤的，就是隐患。跨分
割、换层无伴随地孔、连接器地针不足——带着这三个答
案去审一块板，能抓住大半信号完整性隐患，成本只是几
分钟的追问。

\topic{同时开关噪声 SSN}前面讲的都是线对线的骚扰；
同一片芯片内部的输出也会联手骚扰自己。8 位或 16 位
总线同拍翻转时，每路输出都要给各自的负载电容充放电，
瞬态电流全部经过同一对（或同一小段）电源脚与地脚。
封装引脚、键合丝、板上过孔串起来有几 nH，电流突变
在这点电感上产生 $V=L\,\mathrm{d}I/\mathrm{d}t$：
放电电流涌向地脚，芯片内部的“地”被垫高（地弹，
ground bounce）；充电电流从电源脚灌入，内部 Vcc 被
拉低。芯片内外看到的电源在边沿瞬间塌了一块，塌多少
与同时翻转的路数成正比。

代入数字。8 路输出、每路负载 50pF、摆幅 5V、边沿
5ns：每路瞬态电流
$I=C\,\mathrm{d}V/\mathrm{d}t=50\text{pF}\times5\text{V}/5\text{ns}=50$mA，
8 路合计 400mA，在 5ns 内全部进出同一个地脚。地脚
回路电感取 5nH（一只引脚加过孔的量级），地弹幅度
$V=5\text{nH}\times400\text{mA}/5\text{ns}=0.4$V；
16 路就是 0.8V。后果有两层：这一拍没翻转的输出，其
低电平被垫高 0.4--0.8V——74HC 在 5V 下低电平噪声
容限约 1.2V，被吃掉大半；更隐蔽的是芯片自己的时钟
脚、复位脚若正经过阈值附近，一次总线大翻转就能造出
一个假沿。SSN 的典型症状是“数据总线大动作时，不相
干的地方出错”。早期双列直插封装把电源脚放在对角、
引线又长，SSN 尤其严重；现代封装把电源与地脚成对放
在芯片近旁、多脚并联，正是被这道题逼出来的形状。

缓解 SSN 的每条手段都对应公式里的一个因子：

\begin{longtable}{L{0.2\linewidth}L{0.31\linewidth}L{0.33\linewidth}}
\toprule
手段 & 改变的因子 & 效果量级 \\
\midrule
错开翻转 & 同时翻转路数 $N$ & 8 拆成两拍各 4，噪声减半 \\
增加电源、地脚 & 回路电感 $L$（并联） & 脚位加倍，$L$ 近似减半 \\
减小负载电容 & 每路电流 $C\,\mathrm{d}V/\mathrm{d}t$ & $C$ 减半则电流减半 \\
放慢边沿、调低驱动 & $\mathrm{d}t$ & 边沿加倍，噪声减半 \\
就近去耦 & 电流路径长度 & 瞬态电流不经过引脚电感 \\
\bottomrule
\end{longtable}

执行这张表在各层各有对应动作。芯片厂在现代 FPGA 的
高速 Bank 旁按厂家给定的信号地比配脚、片上集成大量
去耦；板级设计者把总线收发器分组使能、把同拍翻转的
路数拆开，电源入口去耦按 §一的分层做足；写状态机的
人也有份：让同拍变化的输出位数最小的编码，是免费的
SSN 抑制。测量用短接地针的探头（长地线会引入假信号，
见 §五），直接量芯片一只“静止”输出的低电平：总线
翻转那一拍它被垫起多少，就是地弹。鉴别 SSN 与普通
串扰靠翻转路数可变：同一组线上 1 路翻转与 8 路翻转
各跑一次，噪声幅度与路数成正比的是 SSN，与某条特定
邻线的有无强相关的才是串扰。先分清病因，再选上面表
里的手段。

\topic{差分信号为什么抗干扰}差分把一根信号线换成紧
挨着的一对：驱动器输出等大反相的两半各走一条，接收
端只放大两者之差。设两条线各收到同样的干扰 $n$——
远处来的场、参考面的晃动、电源噪声——单端信号里
$n$ 原样叠加在信号上；差分里接收端算
$(s+n)-(-s+n)=2s$，干扰在减法里相消，这就是共模抑
制。相消的前提是“同向等量”：干扰源离两条线越等距、
两条线经历越一致，残留越少——不对称只有 1\% 量级
时，残留就是干扰的 1\%，即 40dB 的抑制。双绞线把两
条线拧在一起、板上把对内两根紧耦合等长布线，都是在
维护这个前提。

把共模噪声的叠加与取差相消画成波形，这个过程一眼可见。

\begin{waveform}[x=0.86cm,y=0.9cm]
  % 共模噪声区间参考虚线
  \draw[BookMuted, dashed, line width=0.5pt] (5.5,2.3) -- (5.5,6.35);
  \draw[BookMuted, dashed, line width=0.5pt] (6.5,2.3) -- (6.5,6.35);
  \node[hw label] at (6.0,6.55) {共模噪声同向叠加};
  % D+
  \draw[BookAccent, line width=0.9pt] (0,4.8) -- (2,4.8) -- (2,5.8) -- (5,5.8) -- (5,4.8) -- (5.5,4.8) -- (5.5,5.15) -- (6.5,5.15) -- (6.5,4.8) -- (7,4.8) -- (7,5.8) -- (9,5.8) -- (9,4.8) -- (10.5,4.8);
  \node[hw label, anchor=east] at (-0.2,5.3) {D$+$};
  % D-（反相，噪声同向）
  \draw[BookAccent, line width=0.9pt] (0,3.6) -- (2,3.6) -- (2,2.6) -- (5,2.6) -- (5,3.6) -- (5.5,3.6) -- (5.5,3.95) -- (6.5,3.95) -- (6.5,3.6) -- (7,3.6) -- (7,2.6) -- (9,2.6) -- (9,3.6) -- (10.5,3.6);
  \node[hw label, anchor=east] at (-0.2,3.1) {D$-$};
  % 差分结果
  \draw[BookTeal, line width=0.9pt] (0,0.3) -- (2,0.3) -- (2,1.5) -- (5,1.5) -- (5,0.3) -- (7,0.3) -- (7,1.5) -- (9,1.5) -- (9,0.3) -- (10.5,0.3);
  \node[hw label, anchor=east] at (-0.2,0.9) {D$+$ $-$ D$-$};
  \node[hw label] at (6.0,-0.25) {取差后噪声相消，波形干净};
\end{waveform}
\diagramnote{共模抑制的过程：噪声 $n$ 同向叠加到 D$+$ 与 D$-$ 上，接收端只放大两者之差，$(s+n)-(-s+n)=2s$，噪声在减法中相消。相消的前提是两条线对称——干扰到两条线的耦合越相等，残留越少。}

回流路径上差分同样占便宜。单端信号的回流必须走参考
面，环路面积由“线到面”的环决定；紧耦合差分对里，
一条线的回流大部分由另一条线承担——去程与回程在对
内部闭合，环路面积缩到对内间距量级，对外辐射和捡拾
外界干扰的能力同步下降。这就是差分对参考面的依赖远
低于单端的原因；不对称的残余仍以共模电流形式存在，
所以参考面依然要，只是要求可以放宽。

小摆幅是第三个红利。LVDS 在 100$\Omega$ 差分端接上
只摆 $\pm350$mV，驱动电流约 3.5mA 且大小恒定、只换
方向，$\mathrm{d}I/\mathrm{d}t$ 天然温和；接收端有
几十到一百 mV 的差值即可可靠判定。摆幅小、电流稳、
环路面积极小，合起来就是低功耗加低辐射。凡是要出板、
出箱、速率上台阶的信号都选了这条路：USB 的 D+ 与
D$-$ 是 90$\Omega$ 差分，PCIe 每个 lane 一对
85--100$\Omega$，LVDS、SATA、以太网同理——不是单
端完全不能用，而是到了那个速率，单端的串扰、辐射与
回流问题一起长大，差分是更省的方案。

差分的全部好处押在“对称”二字上，布线纪律因此很具
体：对内等长——1mm 的长度差就是约 7ps 的相位偏移，
直接把一部分差分信号转成共模，Gbps 级链路要计较到
0.1mm 量级；对内等距、同层、过孔成对；出板之后用双
绞或屏蔽对。教学机上没有真正的差分链路，但有一条可
以直接搬走的直觉：让去的路和回的路贴在一起、经历相
同，干扰就同来同去——单端世界里“信号贴着回流走”
那条规则，与差分本是同一个物理思想。

\topic{EMC 入门：辐射来自大环路与快边沿}电磁兼容
（EMC）问的是最后一个尺度的问题：板子作为整体，向
外辐射多少、能忍受多少。一个电流环就是一圈小天线：
环尺寸远小于波长时，辐射场强
$E\propto I\cdot A\cdot f^2/r$——正比于电流、环路
面积、频率的平方，随距离反比下降。频率平方是最无情
的因子：频率加倍，场强四倍、功率密度十六倍。板的辐
射因此由两个旋钮控制——环路面积与频率成分；本节前
面四段其实都在拧这两个旋钮，发射与敏感还遵守互易：
同一个大环，既善于发射也善于接收。

数字信号的高频含量由边沿而非时钟频率决定：拐点频率
$f_k\approx0.5/t_r$，高于它的谐波迅速衰减。74HC 边
沿约 6ns，$f_k\approx83$MHz；1ns 边沿对应 500MHz；
0.1ns 对应 5GHz。“总线只跑 1MHz”与高频无关是错觉
——边沿决定高频含量，时钟频率只决定能量在谱线上怎
么排。另一个尺度是天线效率：导体长度接近四分之一波
长才是有效天线，10cm 走线对应 $\lambda/4$ 的频率约
750MHz，所以 83MHz 的成分在 10cm 走线上辐射效率其
实很低——这是运气，不是设计。

减缓手段对着两个旋钮加一个盖子。减小环路面积：完整
地平面、回流贴信号走、连接器多插地针——本节前三段
的内容在 EMC 语境下原样重读一遍，就是发射控制。减
缓边沿：用满足速度要求的最慢器件，源端串电阻、调低
驱动强度，把 $f_k$ 从几百 MHz 拉回几十 MHz，辐射按
平方退潮；这条几乎不花钱，只花“不用用不上的速度”
的自觉。屏蔽与接地是盖子：金属机壳接保护地，线缆进
出机壳处加滤波或共模扼流圈，机壳孔缝远小于波长（经
验线取 $\lambda/20$，1GHz 对应约 1.5cm）。屏蔽应当
最后上场：前两条做到位的板子，常常不需要它。排查辐
射超标用的是近场探头：一只小磁环探头贴着板面扫，哪
里场强异常哪里就有大环或快边沿——它量到的正是
$I\cdot A\cdot f^2$ 的分布图。

EMC 还有辐射之外的另一半：传导。30MHz 以下的干扰主
要沿导线进出——电源线、信号线、地线都是通道，治法
是滤波与去耦：电源入口的 $\pi$ 型滤波、线缆出口的
共模扼流圈（差模电流在磁环里磁场相消、畅行无阻，共
模电流磁场相加、遇到高阻），把噪声拦在设备边界上。
30MHz 以上导线越来越像天线，辐射接过主导权，治法回
到环路面积与边沿速度。两者分界并不绝对，但“低频沿
线走、高频向空发”这张地图，足够指导排查的先后：先
看线，再看环。

低速教学机为什么天然友好？把三个因子逐项代入就明白：
时钟 1--10Hz、有效翻转率极低；负载电流 mA 级；74HC
虽有约 80MHz 的边沿成分，但几到十几厘米的走线相对
波长太短，辐射效率很低。$I$、$A$、$f$ 三项里两项半
在低位，面包板不做任何 EMC 措施也不会骚扰谁——它
不是超越了规律，而是无意中符合了规律：电流小、环小、
边沿相对线长而言慢。

\begin{longtable}{L{0.24\linewidth}L{0.27\linewidth}L{0.33\linewidth}}
\toprule
辐射因子 & 面包板教学机 & 高速板 \\
\midrule
时钟与翻转率 & 1--10Hz，极低 & 数十至数百 MHz \\
边沿时间 $t_r$ & 5--8ns（74HC） & 0.1--1ns \\
拐点频率 $0.5/t_r$ & 约 60--100MHz & 0.5--5GHz \\
走线长度与天线效率 & 数厘米，远低于 $\lambda/4$ & 可达数十厘米 \\
综合判断 & 天然低风险 & 必须按本节手段逐项设计 \\
\bottomrule
\end{longtable}

但同样的布线习惯搬到 50MHz 时钟、0.5ns 边沿的板子
上，辐射会按平方规律讨回来。教学机真正可以带走的遗
产是那套检查顺序：先问环多大，再问边沿多快，最后问
回流从哪走——三个问题在任何速率下都成立，这正是本
章把信号完整性放在工程边界来谈的用意：它不是高速板
才有的新问题，而是低速下被宽限、高速下必还清的老约
定。

\section{五、测量与工程实践}

前面四节把电源、时钟、信号完整性各自的机理和算法讲完了，这一节换一个身份：从设计者换成测量者。
测量不是中立的旁观——仪器本身是被测电路的一部分：探头会改变电路，带宽会改写波形，地线会凭空造出振铃。
本节先把两件主力仪器的能力边界算清楚，再讲电源纹波和眼图这两类专项测量，然后用三则调试案例把前四节的规则串起来，最后回答一个经常被问的问题：面包板上这些"宽松"的经验，到了高速世界还剩什么。

一个贯穿本节的观点：每次测量都是两个系统的相遇——被测电路是一个，仪器与探头是另一个。
波形上的每个特征，理论上都要先回答"这是谁产生的"，才有资格回答"这说明什么"。
前四节给的算法正是回答这两个问题的工具：它们既能算电路，也能算仪器；本节反复出现的"自检""换地针复测""改采样率复测"，本质上都是在把两个系统的贡献分离开。

\topic{示波器带宽与探头决定你能看见什么}示波器的前端是一台低通滤波器：超过带宽的频率分量被衰减，到达屏幕时波形已经变形。
物理图景很清楚——方波可以分解成基波加一串奇次谐波，边沿越陡，承担边沿的谐波次数越高；示波器带宽决定了哪些谐波能活着到达屏幕。
工程上用一个经验式把带宽换算成时间：仪器自身可分辨的上升时间 $t_r \approx 0.35 / f_{BW}$。
100MHz 示波器的自身上升时间约 3.5ns，200MHz 约 1.75ns，1GHz 约 0.35ns。

屏幕上显示的上升时间是信号与仪器两者的合成，按方和根相加：$t_{disp} \approx \sqrt{t_{sig}^2 + t_{scope}^2}$。
拿 100MHz 示波器看一个 1ns 边沿：$\sqrt{1^2 + 3.5^2} \approx 3.6$\,ns——屏幕上是一条约 3.6ns 的圆角斜坡，原本可能存在的过冲和振铃（它们由更高的谐波构成）被滤得干干净净。
换句话说，你看到的不是信号，而是信号经过示波器之后剩下的部分；边沿上最快的那些细节，恰好最先被仪器扔掉。
由此得一条选择规则：要"如实看"一个边沿，示波器上升时间应快于信号 3 倍以上（即带宽约 $3 \times 0.35/t_{sig}$），此时显示误差约 5\%；只要求 10\% 误差，仪器上升时间也要快于信号一半左右。

0.35 这个系数值得拆开推一遍，因为它示范了本章的方法论。把示波器前端近似成一阶 RC 低通：阶跃响应按 $1-e^{-t/\tau}$ 爬升，从 10\% 到 90\% 的上升时间为 $\tau\ln 9 \approx 2.2\tau$；同一电路的 $-3$\,dB 带宽为 $f_{BW} = 1/(2\pi\tau)$。
两式相乘消去 $\tau$：$t_r \cdot f_{BW} = 2.2/(2\pi) \approx 0.35$——"带宽"这个频域指标被翻译成了"上升时间"这个时域指标。
真实示波器不止一个极点，系数会略偏离 0.35，但做数量级估算完全够用；遇到任何只给带宽不给时域指标的仪器，都可以先乘一遍 0.35 再说。

换一个等价的频域看法：10MHz 方波由 10MHz 基波加 30MHz、50MHz、70MHz……的奇次谐波叠加而成，边沿越陡，承担边沿细节的谐波次数越高。
经验上要看到"像方波"的形状至少得保留 5 次谐波，要看清 1ns 级边沿，谐波要收到 GHz 门口。
所以 100MHz 示波器看 10MHz 方波只是"够用"：显示出来的圆滑，一半是仪器给的，心里要有数。

还要注意示波器与探头的带宽是各自独立的低通，系统带宽同样按方和根合成：100MHz 示波器配 100MHz 探头，系统带宽只剩约 $100/\sqrt{2} \approx 70$\,MHz。
买仪器时探头标称带宽要留出一档余量，原因就在这里。

把前面的算法串成工作流程：拿到一个测量任务，先问被测信号最快的边沿是多少——查数据手册的输出跳变时间，而不是看时钟频率，决定带宽需求的是边沿；再由 $t_{sig}$ 反推系统带宽下限，核示波器与探头合成后的系统带宽够不够；不够就别测时域形状，改测能量与频谱，或者干脆承认"这台仪器只能看趋势"。
最常见的错误是拿时钟频率当带宽依据：1MHz 的时钟照样可以有 5ns 的边沿，而看住 5ns 边沿需要的带宽是 1MHz 基波的一百多倍。

比带宽更隐蔽的是探头。10x 无源探头标配的鳄鱼夹地线长约 10cm，这段导线加上探头尖端回地夹所围的环路，电感量级约 100--300nH（直导线每厘米约 10nH 上下）。
探头输入电容约 10pF，与这段电感构成一个 LC 回路，谐振频率 $f = 1/(2\pi\sqrt{LC})$；代入 150nH 与 10pF，得约 130MHz——正好落在快边沿的谐波密集区。
于是只要被测边沿足够快，探头回路自己就被激励起来，屏幕上出现一百多 MHz 的衰减振铃；这振铃不是电路的，是测量装置的。
判别办法很便宜：把长地线换成套在探针尖上的短弹簧地针（环路电感降到 10nH 量级），重新测同一点——振铃若明显缩小甚至消失，之前看到的就是测量伪影；振铃若原样不动，才轮到怀疑电路本身。
"同一信号、长地线与短地针各测一次，两张波形并排存下来"是每个实验室都该亲手做一次的实验，比任何说教都有说服力。

10x 探头相对 1x 探头的价值也在这里。1x 探头就是一根屏蔽线，输入电容高达 50--150pF，带宽往往只有十几 MHz，挂到 74HC 的输出脚上就是一个可观的负载，边沿被它自己拖慢。
10x 探头内置 9M$\Omega$ 电阻，与示波器的 1M$\Omega$ 输入阻抗组成十分压：输入阻抗提高到 10M$\Omega$，等效输入电容降到 10pF 上下，带宽可达数百 MHz；代价是信号被衰减 10 倍，示波器底噪相对变大，毫伏级小信号不好测。
测数字边沿、时钟、串扰毛刺时一律用 10x；只有测电源纹波的精细幅度、带宽无所谓时，才考虑 1x。

用 10x 探头前有一道常被跳过的手续：补偿。探头电缆的分布电容与示波器输入电容，会和那 9M$\Omega$ 电阻形成随频率变化的分压误差，探头柄上的小可调电容就是用来把它调平的。
做法是把探头接到示波器面板的 1kHz 校准方波上：波形平顶，补偿正好；前沿冲起尖角，过补偿；前沿圆钝，欠补偿。
跳过这一步，幅度读数可能差出一成，快边沿的形状更是直接失真——很多人测到的"振铃"，源头其实是一颗没调好的补偿电容。

即使是补好的 10x 探头，负载效应也要先估再测。10pF 在 100MHz 下的容抗为 $1/(2\pi f C) \approx 160\,\Omega$——对一个 50$\Omega$ 系统的节点来说，探头一搭上去，等于并联了一个同量级的负载，被测波形当场就变了；1x 探头的 100pF 在同一频率下只有约 16$\Omega$，近乎短路。
估算顺序应该是：先按被测点的阻抗与信号最高频率算出探头引入的负载，再决定测不测、怎么测。
测高速节点时有源探头（输入电容 1pF 上下）存在的意义正在于此；本书的面包板信号源阻抗高、边沿慢，10x 无源探头已经足够宽绰。

还有一个安全层面的常识：示波器的地夹接机壳保护地。测量不以机壳地为参考的电路（开关电源原边、市电相关节点）时，地夹一夹上去就是短路——轻则跳闸烧探头，重则伤人。
有人为此剪断示波器电源线的地脚让整机"浮地"，这是把危险从电路搬到机壳上，坚决不可取；正确做法是差分探头、隔离变压器，或者用两通道分别测两点再做数学相减。
本书全部实验都是低压直流共地系统，不会遇到这个问题，但第一天就该知道这条线划在哪。

带宽与探头之外，显示模式也在做选择。单次触发把一次瞬态定格，适合抓跌落与毛刺；长余辉把成千上万次扫描叠在一起，适合看抖动与偶发异常；平均模式把多次触发的波形逐点平均，随机噪声被压下去、周期性细节浮出来，但它对触发之间会漂移的事件同样无情——平均十次，偶发毛刺就永远消失了。
选模式之前先想清楚要看的是"典型"还是"最坏"：两者都重要，但没有任何一个模式同时擅长。

\begin{longtable}{L{0.16\linewidth}L{0.22\linewidth}L{0.22\linewidth}L{0.28\linewidth}}
\toprule
探头方式 & 输入电容 / 带宽量级 & 典型用途 & 常见误用 \\
\midrule
1x 探头 & 50--150pF / 十几 MHz & 毫伏级小信号、低频模拟 & 拿它看快边沿，负载电容把边沿拖慢还怪电路 \\
10x + 长地线 & 约 10pF / 数百 MHz & 日常数字信号粗看 & 长地线环路自己振铃，把伪影当电路振铃 \\
10x + 弹簧地针 & 约 10pF，环路电感最小 & 边沿、振铃、纹波的定量测量 & 地针没顶在最近的地点上，环路仍然偏大 \\
\bottomrule
\end{longtable}

\topic{逻辑分析仪与时序分辨率}逻辑分析仪与示波器是两种哲学：示波器把少量通道的电压波形尽量测准，逻辑分析仪把几十根线在每个采样点上只判成 0 或 1，丢掉全部模拟细节，换来通道数和存储深度。
代价写在采样率里：100MSa/s 意味着 10ns 才采一个点，比 10ns 窄的毛刺可能从两个采样点之间溜走，或者只在某一个采样点上闪现成一根孤立的细线，真假难辨。
不少机型带毛刺检测（把采样点之间的跳变也记入），但可捕获的最窄脉宽仍由专门的毛刺时钟决定；读规格书时找"最小可捕获脉宽"这一项，它比宣传的最大采样率更接近真相。

采样还带来一个与示波器共有的陷阱：混叠。采样定理要求采样率高于信号最高变化频率的两倍，否则高频成分被折叠成低频假象——用 50MSa/s 去采一个 48MHz 的时钟，屏幕上会显示一个 2MHz 的"信号"，而电路里根本不存在这样的东西。
判别办法是改采样率再测一次：真信号的频率不随采样率动，混叠产物会跟着跑。
所以定时模式的采样率一般取被测最快信号的 4--10 倍，既躲混叠，也留足看清先后关系的余量。

逻辑分析仪的另一半价值在触发与存储。触发条件可以写成码型（"地址为 \code{0xF0} 且写使能有效"）、序列（"先看到 A 再看到 B"）或毛刺；仪器在环形缓冲里不停地采，条件命中时把命中点前后各一段定格——偶发故障因此第一次变得可以反复研究。
存储深度与采样率互相抢资源：观察窗口 = 深度 ÷ 采样率，1M 点深度配 100MSa/s 只有 10ms 窗口；想看更长的时间跨度就得降采样率，这正是"看得清"与"看得久"的经典交换，买机器和设参数时是同一笔账。
新机型的协议解码（UART、I$^2$C、SPI 等）把逐拍的 0/1 直接翻译成帧，调串行总线效率倍增；但要记住解码建立在采样判决之上，模拟层的问题——缓慢边沿、擦着阈值的毛刺——会被它悄悄抹平，协议层怎么也说不通时，仍要回到示波器看模拟波形。

接线也有纪律：每路信号的参考地要就近接，探头飞线越短越好——逻辑分析仪探头同样受本章全部寄生参数支配，几十根飞线捆成一束，本身就是一个串扰丛生的结构。
阈值电平要按被测逻辑家族设定：5V CMOS 与 3.3V LVTTL 的判决点不同，设错了，0 和 1 的边界本身就是错的，后面全部结论跟着作废。

举个实战例子把触发写具体：怀疑第七章的总线在某拍被两个源同时驱动，就设码型触发——"时钟上升沿且 \code{CO}=1 且 \code{RO}=1"，命中一次就抓到一次争用的瞬间；再设序列触发"LDA 的 T3 之后紧跟 STA 的 T2"，可以核对指令切换的边界拍。
逻辑分析仪与示波器还能联合作战：很多机型有触发输出端子，条件命中时吐出一个脉冲，把它接进示波器的外部触发，示波器就在"逻辑分析仪认定的那个时刻"展开模拟波形——时序定位与波形细节同时对齐，偶发故障的最后一层隐身术就此解除。

由此有分工：查多信号之间的相对时序——8 位总线、十几根控制线、地址有效到数据有效的先后——用逻辑分析仪；查单根线上的波形细节——过冲多高、振铃什么频率、毛刺多宽——用示波器。
逻辑分析仪还有一种与被测电路对齐的用法：状态模式下用电路自己的时钟作采样时钟，每个时钟沿记录一次全部通道，得到一张逐拍的状态表，与第七章的微码表逐行对照最方便。
这等于把"仪器视图"对齐到"设计视图"：设计是按时钟拍的，状态模式看到的就是时钟拍的。

教学实验里怎么选：调第七章的微码序列、核对总线争用、追一条控制线哪一拍被意外置位，逻辑分析仪一次看全 8 根总线加全部控制线，是主力；查电源纹波、时钟振铃、复位毛刺，通道只要两三个但波形细节就是全部信息，用示波器。
一个粗糙的判别口诀：要问"哪一拍、谁先谁后"，用逻辑分析仪；要问"多高、多快、什么形状"，用示波器。

抓到的数据别浪费：逻辑分析仪的记录可以导出成逐拍表格，与微码仿真器的 trace 逐行对拍——差异出现的第一行，就是硬件与设计的第一处分歧。
这一步把调试从"盯着屏幕找不同"变成可重复的对照实验，也是第七章逐拍 trace 方法论在仪器上的延伸。

\begin{longtable}{L{0.20\linewidth}L{0.34\linewidth}L{0.34\linewidth}}
\toprule
维度 & 示波器 & 逻辑分析仪 \\
\midrule
通道数量级 & 2--4 个 & 16--34 个起步 \\
看到的东西 & 电压随时间的连续形状 & 每个采样点上的 0/1 判决 \\
分辨率由什么定 & 带宽与上升时间 & 采样率与最小可捕获脉宽 \\
典型任务 & 振铃、过冲、纹波、边沿时间 & 多线相对时序、协议、逐拍状态 \\
本书实验里的用法 & 测纹波、查毛刺、验时钟质量 & 核对微码序列、抓总线争用 \\
\bottomrule
\end{longtable}

\topic{电源纹波怎么测}电源测量的第一个敌人是直流本身：5V 上叠着 50mV 纹波，DC 耦合下整条波形被顶在屏幕上方，纹波只占垂直方向的百分之一，什么都看不清。
把通道切到 AC 耦合，信号路径里串入一个隔直电容，直流被挡住，只剩变化部分，于是可以用 10--20mV/div 的档位把纹波展开。
第二个敌人是拾取：开关电源、时钟线、继电器都在周围辐射，探头地线围成的环路就是一副小天线。
纹波测量要用最短的地——弹簧地针直接点在去耦电容两端，而不是夹着一根长地线去够远处的电源脚；示波器若有 20MHz 带宽限制，打开它，挡住与电源无关的高频拾取。

探测点的选择同样要先想清楚：纹波要测在负载端——目标芯片电源脚旁的去耦电容两端，那才是芯片真实看到的电源；测在稳压器输出端，读到的是供电路径的起点，中间走线的压降与跌落全被漏掉。
排查时可以沿供电路径逐点测量：从稳压器输出、过板级电容、到芯片引脚，跌落深度在哪里突变，路径阻抗就藏在哪里。

测到波形之后要会分类，因为两类波形指向两类病根。
稳态纹波是周期性的：线性电源之后是工频整流留下的 100Hz 锯齿，开关电源之后是开关频率（几十 kHz 到几 MHz）的三角波或尖刺，幅度基本不随负载突变。
瞬态跌落是事件性的：负载阶跃（电机启动、继电器吸合、一片芯片几十个输出同时翻转）之后电压下陷再回升，持续几 µs 到几百 µs，要用单次触发或长余辉去抓。
稳态纹波大，问题多在稳压器与滤波；瞬态跌落深，问题多在去耦储能不足或供电路径阻抗大——第一节的目标阻抗算法管的就是后者。

两类波形都能再拆细、再演算。线性稳压器输出的稳态纹波，主要是整流滤波后残余的 100Hz 波动穿过稳压环路：LDO 的电源抑制比在低频有 60--70dB，但随频率升高迅速变差，到几百 kHz 往往只剩 20--30dB，输入上的高频波动会被大量放行——指望 LDO 挡住开关预稳压器的全部噪声，是不读曲线的想当然。
开关电源的纹波由两部分组成：电感电流纹波在输出电容 ESR 上产生的三角波（50m$\Omega$ 的 ESR 配 200mA 峰峰值纹波电流，就是 10mV 峰峰值），加上开关管动作瞬间的振铃尖刺——后者能量小但频谱很高，正是 20MHz 带宽限制要挡住的东西。
瞬态跌落同样先估再测：100mA 的负载阶跃全靠一颗 10µF 本地电容支撑 10µs，跌落为 $\Delta V = I\Delta t/C = 0.1 \times 10/10 = 0.1$\,V；够不够、该由哪一层补，看跌落深度与持续时间的比例就明白。
测到的波形因此不是一堆"噪声"，而是一条可以反推 PDN 在哪个时间尺度上缺电容的诊断曲线——这与第一节按时间尺度分层的思路首尾相接。

读数也有讲究：纹波按峰峰值读，不按有效值——决定逻辑是否误翻转的是最坏的瞬间，不是平均能量。
测量时探头方向、地的接点、带宽限制状态都要写进记录，因为换任何一项读数都会变；两次测量差出 20\%，往往不是电源变了，而是测法变了。
有条件的实验室还可以把示波器的 FFT 打开，或直接上频谱仪：稳态纹波在频谱上是一根根谱线，谱线位置直接指认来源——100Hz 是整流残余，开关频率及其谐波是电源本身，与时钟频率重合的那根则要怀疑耦合而不是供电。
配一个近场小环探头在板子上扫一遍，辐射最强的位置往往就是回流路径最拧巴的位置，这是把第四节 EMC 机理变成定位手段的用法。

最常见的误判，是把探头环路拾到的磁场感应当成电源噪声。
判别方法叫双点同地：把探头尖端与地针同时点在同一个地焊盘上，理论上应显示一条平线；若此时屏幕上仍有"纹波"，那它是环路拾取的干扰，不是电源的。
先做这个自检再下结论，是电源测量的纪律；整个流程可以写成一张短清单，每次照做：

\begin{lstlisting}
纹波测量流程
1. 10x 探头换弹簧地针, 点在目标芯片的去耦电容两端
2. AC 耦合, 10-20mV/div, 打开 20MHz 带宽限制
3. 双点同地自检: 尖与地针同点一处, 平线才算环境干净
4. 看稳态纹波: 频率、峰峰值, 与开关频率对照
5. 单次触发抓瞬态: 用负载动作信号做外触发
6. 记录: 波形 + 探头接法, 缺一不复现
\end{lstlisting}

\topic{眼图：把时序裕量画成一只眼睛}让示波器用远高于比特率的采样率连续采集一段数据波形，再按一个比特周期（UI，单位间隔）把波形切成段、全部叠到同一张图上：所有 0→1、1→0、0→0、1→1 的轨迹交织在一起，中央留出一块空白，整张图形像一只睁开的眼睛，这就是眼图。
眼睛的张开程度把两个抽象裕量变成一眼可见的形状：垂直方向张开多高，是判决时刻的噪声裕量；水平方向张开多宽，是抖动留给采样点的时序裕量；交叉点那团模糊带的宽度，就是抖动的直接读数。
高速链路的标准干脆用模板（mask）说话：在眼睛中央画一个禁止进入的多边形，发射端在所有码型下都不许让波形碰它；接收端依此设计判决门限与采样时刻。

本书的面包板用不上眼图：时钟周期以 µs 计，裕量以百 ns 计，眼睛永远全开，测它没有信息量。
但要理解千兆位链路为什么人人谈眼图：1Gbps 的 UI 只有 1ns，5Gbps 只剩 200ps，抖动和噪声各吃掉一大块之后，眼睛只剩一条缝；还能不能可靠判决，先看眼图，再看误码率。
眼图的本质是"把统计学画成一张图"：它叠的是成千上万个周期，显示的是最坏情况出现的频率，而不只是一次偶然的波形——这正是单次触发波形给不出的视角。

读眼图先认部位：中央开口的高度叫眼高，对应噪声裕量；开口宽度叫眼宽，对应时序裕量；上下两条"眼皮"的厚度，是幅度噪声与反射的堆积；交叉点的横向模糊带是抖动。
规范做法是统计判决时刻处 0 电平与 1 电平各自分布的交叠程度——交叠越深，误码越近。
抖动本身也要分两类看：确定性抖动有界（来自码型相关、串扰、占空比失真），随机抖动无界、服从高斯分布（来自热噪声与时钟源的相位噪声）。
把误码率对采样时刻画成曲线，就是著名的"浴缸"：中间一段误码率低到测不出，两端沿随机抖动的高斯尾部陡升；两壁间距再扣掉规范要求的置信水平，才是链路真正能用的时序裕量。
眼图是浴缸曲线的直观版——前者给人看，后者给标准算。

把各部位标在一张示意眼图上，对照着读更直观。

\begin{pdfdiagram}
  % 叠加的码元轨迹（成对绘制，形成眼皮厚度）
  \draw[BookAccent, line width=0.85pt] (0,3.1) -- (0.3,3.1) -- (1.3,0.9) -- (8.7,0.9) -- (9.7,3.1) -- (10,3.1);
  \draw[BookAccent, line width=0.85pt] (0,0.9) -- (0.3,0.9) -- (1.3,3.1) -- (8.7,3.1) -- (9.7,0.9) -- (10,0.9);
  \draw[BookAccent!70, line width=0.7pt] (0,2.92) -- (0.45,2.92) -- (1.55,0.72) -- (8.85,0.72) -- (9.85,2.92) -- (10,2.92);
  \draw[BookAccent!70, line width=0.7pt] (0,0.72) -- (0.45,0.72) -- (1.55,2.92) -- (8.85,2.92) -- (9.85,0.72) -- (10,0.72);
  % 模板六边形（禁止进入区）
  \draw[BookRed, dashed, line width=0.7pt] (4.0,2.0) -- (4.6,2.5) -- (5.4,2.5) -- (6.0,2.0) -- (5.4,1.5) -- (4.6,1.5) -- cycle;
  \node[hw label, text=BookRed] at (5,2.0) {模板};
  % 眼高
  \draw[{Stealth[length=1.8mm]}-{Stealth[length=1.8mm]}, BookMuted, line width=0.6pt] (6.5,0.9) -- (6.5,2.92);
  \node[hw label, anchor=west] at (6.65,1.9) {眼高};
  % 眼宽
  \draw[{Stealth[length=1.8mm]}-{Stealth[length=1.8mm]}, BookMuted, line width=0.6pt] (1.55,1.22) -- (8.45,1.22);
  \node[hw label] at (5,1.08) {眼宽};
  % 交叉点
  \node[hw label] at (1.7,0.32) {交叉点模糊带 = 抖动};
  % UI 标尺
  \draw[BookMuted, line width=0.6pt] (0.8,-0.15) -- (9.2,-0.15);
  \draw[BookMuted, line width=0.6pt] (0.8,-0.25) -- (0.8,-0.05);
  \draw[BookMuted, line width=0.6pt] (9.2,-0.25) -- (9.2,-0.05);
  \node[hw label] at (5,-0.42) {一个比特周期（UI）};
\end{pdfdiagram}
\diagramnote{眼图的形成与读法：把数据波形按一个 UI 切段叠加，所有跳变轨迹交织出中央开口。眼高对应判决时刻的噪声裕量，眼宽对应采样点的时序裕量，交叉点的横向模糊带是抖动的直接读数；红色六边形是规范模板，任何轨迹都不得进入。}

\begin{longtable}{L{0.20\linewidth}L{0.28\linewidth}L{0.40\linewidth}}
\toprule
眼图部位 & 对应物理量 & 典型恶化来源 \\
\midrule
眼高（开口高度） & 判决时刻的噪声裕量 & 反射、串扰、信道衰减、电源噪声 \\
眼宽（开口宽度） & 采样点的时序裕量 & 抖动、码间干扰、偏斜 \\
眼皮厚度 & 幅度噪声与多重反射的堆积 & 端接不良、电源波动 \\
交叉点模糊带 & 抖动总量的直接读数 & 时钟源相位噪声、串扰、码型相关 \\
交叉比（交叉点电平） & 判决门限的对称程度 & 上下拉驱动不对称、阈值偏移 \\
\bottomrule
\end{longtable}

主流高速标准都把模板测试写成明文：USB、PCIe、DDR 各自规定一张禁止进入的模板和配套测试夹具，发射端在最坏码型、最坏温度下都不许让波形碰模板边界。
模板的意义，是把"链路能不能工作"这个系统问题，拆成发射端、信道、接收端三方可以各自独立验证的指标——这是工程上处理复杂系统的标准手法，与软件里的接口规格同源。

面包板上其实可以做一个概念演示：用函数发生器出伪随机码流，串一级 RC 把带宽故意压窄，示波器用码流的同步时钟触发并开长余辉——眼睛会肉眼可见地闭上，RC 越大闭得越紧。
这个演示不能量化任何裕量，但能让你亲手把"带宽不足→眼图闭合→判决变难"这条因果链走一遍；将来在高速链路的规格书里再遇到眼图模板，看到的就不再是黑话，而是当年那只被 RC 慢慢捏合的眼睛。

把眼图与本章其他内容连起来看：眼高被串扰与电源噪声啃——那是第四、第一节的事；眼宽被抖动与偏斜啃——那是第二节的事；眼皮的形状由反射决定——那是第三节的事。
一张眼图，就是全章算法的一次联合阅卷。

\topic{调试案例三则}以下三则案例分别对应第一、二、四节的机理，按"症状、定位、根因、修复"各讲一遍；排查手段都不超出本节讲过的仪器纪律。

案例一：大电流外设动作时 MCU 随机复位。
\textbf{症状}系统平时稳定，每当继电器吸合（或电机启动、电磁阀动作）的瞬间，MCU 有相当概率复位；外设不动作时连跑几天不出错。
\textbf{定位}示波器 AC 耦合、弹簧地针，探在 MCU 电源脚旁的去耦电容两端，用外设动作信号作外触发：在继电器吸合时刻抓到一次几十 µs 的电压深坑，谷底跌破复位芯片的阈值；坑深与吸合电流大小正相关。
\textbf{根因}继电器线圈几百 mA 的吸合电流与 MCU 共用一段细长电源走线，走线的电阻加寄生电感在瞬态下压出深坑（$L\,di/dt$ 项尤其凶：走线电感到几百 nH、电流在 µs 内变化几百 mA 时，坑深可达伏级）；板级储能电容离外设又远，来不及补给。用第一节的语言说：在这次瞬态的时间尺度上，电源分配网络没有就近的电荷库。
\textbf{修复}外设电源入口处就近并 100--470µF 电解电容作本地储能；外设与 MCU 的供电从电源出口起星形分开走线，不再共用细线段；MCU 一侧维持原有 100nF 加 10µF 的本地去耦。改后同一触发下重测，深坑应缩到复位阈值之上一个安全距离。

案例二：同一片上两寄存器偶尔传错数据。
\textbf{症状}某条寄存器到寄存器的传输绝大多数时候正确，偶发错一个值；温度升高或换一批芯片后出错率变化；把时钟频率降下来就不出错——看起来像建立时间被吃掉了。
\textbf{定位}逻辑分析仪状态模式逐拍记录数据通路，把出错的拍挑出来，发现错的是"新值没被采到、采到的还是旧值"；再用示波器两通道分别点发射端与接收端的时钟脚，测得两个时钟沿之间有约 0.3ns 的固定差距，方向是接收端早到。
\textbf{根因}时钟树两条分支长度差约 5cm（FR4 上信号传播约 60--70ps/cm，5cm 即 0.3ns 量级），其中一支又多经过一级缓冲器，偏斜方向不利；第二节的建立时间预算按"理想零偏斜"留的裕量，实际被偏斜、抖动、温度漂移分掉，最坏组合下为负。
\textbf{修复}重布时钟分支使两支等长，拉平缓冲器级数；来不及改板，就在该路径上减一级组合逻辑或换更快的器件，把数据路径延迟收回预算内。长期办法是把偏斜当作预算表的正式一项，不再假设为零；改后用高温、满速做长时间回归。

案例三：复位线上的时钟毛刺造成偶发复位。
\textbf{症状}系统偶发复位，但没人碰复位按钮、看门狗也没到期；出错与程序行为无关，却在时钟负载最重的外设工作时更频繁。
\textbf{定位}示波器两通道同时看时钟线与复位线：复位线上每逢某些时钟沿就出现一个窄毛刺，宽几 ns 到十几 ns，与时钟沿严格对齐、极性同向；毛刺幅度只在复位输入阈值附近徘徊，所以只有一部分毛刺触发复位——"偶发"由此而来。
\textbf{根因}复位线与时钟线在板上长距离平行走线，按第四节的串扰机理，容性耦合与时钟沿的磁场变化都在邻线上感应电压；平行越长、间距越近，毛刺越大。复位输入又是高阻、无滤波的敏感端，几 ns 的毛刺也认账。
\textbf{修复}能改板：复位线与时钟线拉开间距（至少 3 倍线宽量级）、缩短平行段，必要时中间隔一根接地护线。不能改板：复位输入端加 RC 滤波（如 10k$\Omega$ 加 100nF，时间常数 1ms，远大于毛刺宽度）再经施密特触发器整形，几 ns 的毛刺被积分掉。改后复测：毛刺仍在，但进不了阈值，复位不再误发。

\begin{longtable}{L{0.20\linewidth}L{0.22\linewidth}L{0.26\linewidth}L{0.22\linewidth}}
\toprule
案例 & 关键机理 & 首选定位手段 & 修复方向 \\
\midrule
外设动作致复位 & 瞬态电流在无本地储能的路径上压出深坑 & AC 耦合 + 外触发，抓电压深坑 & 就近储能 + 星形供电 \\
寄存器偶发传错 & 时钟偏斜吃掉建立时间裕量 & 状态模式挑错拍 + 两通道测时钟沿 & 等长分支 + 压缩数据路径延迟 \\
复位线误触发 & 平行走线串扰出窄毛刺 & 两通道看毛刺与时钟沿的对齐 & 拉开间距 + RC 滤波整形 \\
\bottomrule
\end{longtable}

三则案例的共同结构值得再看一眼：症状都是"偶发"，定位都靠"把偶发变成必现"——外触发让深坑每次都被抓到，状态模式让错拍可以被挑出，双通道对齐让毛刺与时钟沿的因果关系无可抵赖。
偶发故障不是不可复现，只是复现条件藏得深；仪器纪律的全部意义，就是把藏着的条件挖出来。
把三则案例的共通动作抽出来，就是一张排查清单，顺序比内容更重要——先让故障变得可重复，再谈根因：

\begin{lstlisting}
偶发故障排查顺序
1. 定义"偶发": 与什么事件相关(外设动作/温度/批次/上电时长)
2. 造触发: 用相关事件做示波器外触发, 或逻辑分析仪码型触发
3. 固定对照: 一次只改一个变量, 保留未出错时的原始记录
4. 先排测量伪影: 换短地针、改采样率, 确认真假再下结论
5. 定位到时间与位置: 哪一拍、哪个引脚、沿什么路径耦合
6. 根因归类: 电源/时钟/串扰, 回到本章对应的预算表
7. 修复后回归: 最坏条件(高温、满速、大负载)长时间复测
\end{lstlisting}

三则之外还有两类常客值得点名。一是接触不良：面包板簧片疲劳、焊点虚焊，症状是"拍一拍板子行为就变"，任何触发都抓不住它，只能靠逐区按压和补焊排查。
二是静电与闩锁：干燥环境里人体带的静电打入输入脚，轻的让芯片内部状态翻转，重的触发寄生可控硅结构、电源电流猛增——症状同样是"偶发"，且与操作员走近的顺序相关。
这两类不归本章任何一张预算表管，但它们共享同一条方法论：先找相关性，再造复现，最后才谈机理。

最后补一条贯穿所有案例的习惯：记测量日志。每次测量写下日期、仪器设置、探头接法、触发条件与结论，波形截图编号存档——"测不到"也要记，它是排除法的一部分。
调试偶发故障常常横跨几天，记忆会骗人，日志不会；而且日志强制你把每次测量的前提写清楚，很多"互相矛盾"的测量结果，一对日志就发现是接法不同。

\begin{lstlisting}
测量日志条目示例
2026-05-11 纹波复测(案例一改板后)
设置: CH1 AC耦合 20mV/div, 20MHz带宽限制, 10x探头
接法: 弹簧地针点在 MCU 电源脚旁 C7 两端(负载端)
触发: 继电器吸合信号外触发, 单次捕获
结果: 跌落 120mV 持续 40µs, 谷底 4.88V
对照: 改板前跌落 320mV, 谷底 4.68V
结论: 谷底已高于复位阈值 4.6V 有余量, 转入长时间回归
\end{lstlisting}

\topic{低速实验与高速世界的关系}本书全部实验生活在一个安全区里：面包板上走线几厘米、时钟周期以 µs 计、74HC 边沿几 ns 到十几 ns。
把数字摆开：5cm 走线的传播延迟约 0.3ns，只是边沿时间的几十分之一，集总模型稳稳成立；瞬态电流以 mA 到百 mA 计，100nF 加 10µF 的两级去耦足够宽绰；探头接法随意一点，波形也大差不差。
安全区里的规则是"建议"，违反了大不出丑。
同一批规则到 100MHz 以上全部变成硬约束：周期 10ns、边沿 1ns 上下时，5cm 走线延迟已达边沿时间的三分之一，走线必须按传输线端接；0.3ns 偏斜占周期 3\%，预算表必须逐项列出；瞬态电流以 A 计、以 ns 计，去耦要按目标阻抗设计；探头与治具不讲究，测量本身就是噪声源。
但请注意：机理一条没变，变的只是数量级把寄生参数从"可忽略"推到了"必须算"。
所以从低速到高速的迁移不是换一批公式，而是换一套检查的严格程度——下表把本书经验逐条对应过去，左边每条你都在面包板上见过，右边每条都是 100MHz 之后必须做的功课。

\begin{longtable}{L{0.16\linewidth}L{0.36\linewidth}L{0.36\linewidth}}
\toprule
主题 & 本书面包板上的经验形态 & 100MHz+ 世界的对应约束 \\
\midrule
走线长度 & 几 cm 飞线，延迟可忽略 & 延迟与边沿可比，按传输线处理并端接 \\
边沿速度 & 74HC 边沿几 ns，慢而安全 & 边沿 1ns 以内，谐波到 GHz，EMC 成问题 \\
地 & 面包板地排栅加星形，大致等电位 & 完整地平面是硬前提，回流贴着信号走 \\
去耦 & 每片 100nF 加每区 10µF，按惯例放 & 按目标阻抗分层设计，仿真与实测闭环 \\
时序 & 裕量以百 ns 计，闭眼都够 & 预算逐项列表，偏斜与抖动各占一行 \\
测量 & 长地线也测得出大概 & 探头即电路的一部分，治具与校准先行 \\
\bottomrule
\end{longtable}

表里的每一行都能还原成一组数字对照。面包板：1MHz 时钟、周期 1µs、边沿约 5--10ns、5cm 走线延迟约 0.3ns——延迟只占边沿的 3\%、周期的 0.03\%，寄生参数全线失守不了。
100MHz 数字板：周期 10ns、边沿约 1ns，同样 5cm 走线的 0.3ns 已占边沿的三分之一、周期的 3\%；谐波伸到 GHz，任何一段没端接的走线都是潜在天线，任何一块被割开的地平面都让回流绕路。
同样的走线、同样的接法，变的只是时钟与边沿，结论就从"怎么都行"翻成"步步要算"——这正是信号完整性作为一门学科存在的理由。

真拿到一块高速任务，检查顺序可以沿用本书的习惯，只是每一项都要量化：先问边沿多快（决定带宽、端接与 EMC 的一切），再问走线多长（决定要不要传输线模型），接着查回流路径（地平面完整吗、跨分割了吗），然后查去耦分层（各时间尺度的电容都在位吗），最后查时序预算（偏斜与抖动逐项列出了吗）。
五问下来，大多数低速经验会原样通过；通不过的那几项，就是要按本章算法重新计算的地方。

什么时候该离开面包板、上正式 PCB？判据可以直接从本章读出：时钟频率上到 10MHz 量级、边沿要求快于几 ns、电源电流上到安培级、或者出现毫伏级模拟小信号——满足任何一条，面包板的寄生参数就从"宽绰"变成"捣乱"。
离开之前，先在面包板上把好习惯养出来；下面这张清单每条都对应本章一处算法，到了 PCB 上只会更严格：

\begin{lstlisting}
面包板阶段就要养成的测量与设计习惯
1. 探头常备弹簧地针, 长地线只用于粗看
2. 任何振铃先换地针复测, 再怀疑电路
3. 电源测量先做双点同地自检, 再读数
4. 时序问题先写预算表, 再上仪器
5. 偶发故障先造触发条件, 再谈根因
6. 每次测量记日志, "测不到"也记
\end{lstlisting}

这套迁移眼光也是本书通往工程实践的桥：你在面包板上养成的习惯——先问物理原因、再估数量级、最后测量复核——正是高速设计每天在做的事，只是那里的估算必须更准、测量必须更讲究。
仪器教会人的最后一件事是谦卑：屏幕上每个波形，都是电路与测量系统共同的作品；分清哪一部分是谁的，是硬件工程师的基本功。

\section*{章末练习}

\begin{longtable}{L{0.2\linewidth}L{0.5\linewidth}L{0.2\linewidth}}
\toprule
任务 & 要求 & 学习目标 \\
\midrule
LDO 功耗演算 & 9V 输入、5V 输出的 LDO 带 200mA 负载：求 LDO 自身功耗与效率；输入改为 12V 后两者各变多少？无散热片的 TO-220 封装热阻约 50°C/W，哪种输入下会过热 & 压差乘电流就是发热 \\
去耦分层选择 & 某芯片输出翻转时在 10ns 内需要 50mA 瞬态电流，电源脚允许跌落 50mV：由 $\Delta Q = C\Delta V$ 求旁路电容下限，并说明它属于哪一层去耦、必须放在哪里 & 瞬态由最近一层供给 \\
偏斜重算预算 & 时钟周期 10ns，两寄存器间数据路径最大延迟 6ns、最小延迟 1ns，建立时间 2ns、保持时间 1ns：偏斜为 0 时建立与保持裕量各多少？若接收端时钟晚到 1.5ns，哪个裕量变好、哪个被击穿 & 保持违例不能靠降频挽救 \\
反射系数计算 & 50$\Omega$ 传输线末端分别开路、短路、接 50$\Omega$、接 150$\Omega$：求四种情形的反射系数，入射阶跃 1V 时反射波各是多少伏、末端电压各是多少 & $\Gamma=(Z_L-Z_0)/(Z_L+Z_0)$ \\
串扰间距规则 & 两条 5cm 平行走线紧挨（间距 0.6mm），时钟沿在邻线上感应出 0.8V 毛刺；按耦合随间距近似反比下降估算，把毛刺降到 0.1V 以下至少要拉开到多少？不能改间距时给出两条替代措施 & 间距、平行长度、护线 \\
示波器带宽选择 & 要如实观察 74HC 输出约 5ns 的上升沿，要求显示的上升时间误差不超过 10\%：示波器带宽至少多少？手头 100MHz 机型够不够，差多少 & $t_r \approx 0.35/f_{BW}$ 与方和根合成 \\
\bottomrule
\end{longtable}

\section*{参考解答与要点}

\begin{longtable}{L{0.18\linewidth}L{0.5\linewidth}L{0.22\linewidth}}
\toprule
任务 & 参考解答 & 必须掌握的要点 \\
\midrule
LDO 功耗演算 & $P = (9-5)\,\mathrm{V} \times 0.2\,\mathrm{A} = 0.8$\,W，效率 $= 5/9 \approx 56\%$；12V 输入时 $P = 7 \times 0.2 = 1.4$\,W，效率降到 $5/12 \approx 42\%$。无散热片 TO-220 按 50°C/W：0.8W 温升约 40°C，室温下勉强可用；1.4W 温升约 70°C，叠加环境温度后越过结温红线，必须加散热片或改用开关电源 & LDO 效率约等于 $V_{out}/V_{in}$，压差全部变成热 \\
去耦分层选择 & $C \geq I\Delta t/\Delta V = 50\,\mathrm{mA} \times 10\,\mathrm{ns} / 50\,\mathrm{mV} = 10$\,nF，取标准 100nF 陶瓷留一个数量级裕量。10ns 时间尺度的电荷只能由最近一层供给：它属于芯片旁去耦层，必须贴在电源脚边、回路总长压到毫米级；局部 1--10µF 与板级 10--100µF 管更慢、更大的瞬态，救不了这一拍 & $\Delta Q = C\Delta V$；时间尺度决定由哪一层供 \\
偏斜重算预算 & 偏斜为 0：建立裕量 $= 10-6-2 = 2$\,ns，保持裕量 $= 1-1 = 0$\,ns，刚好贴边。接收端晚到 1.5ns：建立裕量 $= 10+1.5-6-2 = 3.5$\,ns，变好；保持裕量 $= 1-1.5-1 = -1.5$\,ns，被击穿——发射端下一拍的新数据在接收端采样沿之后 1ns 内就到了。注意保持检查与时钟周期无关，降频救不回来，只能改数据路径延迟或时钟偏斜 & 偏斜对建立与保持的作用方向相反；保持违例与周期无关 \\
反射系数计算 & $\Gamma = (Z_L - Z_0)/(Z_L + Z_0)$。开路 $\Gamma = +1$：反射 +1V，末端 2V；短路 $\Gamma = -1$：反射 $-1$V，末端 0V；接 50$\Omega$ 时 $\Gamma = 0$：无反射，末端 1V；接 150$\Omega$ 时 $\Gamma = 100/200 = +0.5$：反射 +0.5V，末端 1.5V & 阻抗连续处无反射；$\Gamma$ 的符号决定回波极性 \\
串扰间距规则 & 反比估算：间距需扩大 $0.8/0.1 = 8$ 倍，$0.6\,\mathrm{mm} \times 8 \approx 4.8$\,mm，取 5mm 以上。替代措施：两线之间插一根接地护线吸收大部分耦合；缩短平行段长度（耦合与平行长度近似正比）；还可在受害线输入端加 RC 滤波、或给时钟线源端串电阻减缓边沿 & 串扰随间距快速下降、随平行长度正比增长 \\
示波器带宽选择 & 要求 $t_{disp} = \sqrt{t^2 + t_{scope}^2} \leq 1.1 \times 5$\,ns，解得 $t_{scope} \leq 5\sqrt{1.21-1} \approx 2.3$\,ns，带宽 $\geq 0.35/2.3\,\mathrm{ns} \approx 150$\,MHz。100MHz 机型自身上升时间 3.5ns，显示值 $\sqrt{25+12.25} \approx 6.1$\,ns，误差约 22\%，不够；需要 150--200MHz 档位 & 10\% 误差要求仪器上升时间快于信号一半左右 \\
\bottomrule
\end{longtable}
